% From Hartree–Fock to Coupled Cluster: Theory Notes
% Build with the XeTeX engine (fontspec), e.g.  tectonic main.tex
\documentclass[a4paper,11pt,openany]{book}
\usepackage{fontspec}

\usepackage[margin=2.3cm]{geometry}
\usepackage{amsmath,amssymb,amsthm,mathtools,bm}
\usepackage{braket}
\usepackage{booktabs,longtable,array}
\usepackage{xcolor}
\usepackage{listings}
\usepackage[most]{tcolorbox}
\usepackage[colorlinks=true,linkcolor=blue!60!black,citecolor=green!40!black,urlcolor=blue!60!black]{hyperref}
\usepackage{enumitem}


\setlength{\emergencystretch}{3em}

\raggedbottom  % do not stretch pages to the bottom margin: a box that does not fit would leave large gaps
% Notation macros of main.tex and the chapters
\newcommand{\cre}[1]{a^{\dagger}_{#1}}          % creation operator
\newcommand{\ann}[1]{a_{#1}}                    % annihilation operator
\newcommand{\asym}[2]{\langle #1\Vert #2\rangle} % <pq||rs>
\newcommand{\phys}[2]{\langle #1\vert #2\rangle} % <pq|rs>
\newcommand{\chem}[2]{(#1\vert #2)}             % (pq|rs)
\newcommand{\no}[1]{\{#1\}}                     % normal order {...}
\newcommand{\Hbar}{\bar{H}}
\newcommand{\HN}{H_{\mathrm{N}}}
\newcommand{\FN}{F_{\mathrm{N}}}
\newcommand{\VN}{V_{\mathrm{N}}}
\newcommand{\Ehf}{E_{\mathrm{HF}}}
\newcommand{\Ecorr}{E_{\mathrm{corr}}}
\newcommand{\Pij}[1]{P(#1)}                     % antisymmetrizing permutation operator
\newcommand{\tr}{\operatorname{tr}}
\newcommand{\diag}{\operatorname{diag}}
\newcommand{\occ}{\mathrm{occ}}
\newcommand{\vir}{\mathrm{vir}}
\newcommand{\dd}{\,\mathrm{d}}
\newcommand{\code}[1]{\texttt{#1}}
\newcommand{\eqnote}[1]{\quad\text{(#1)}}
\newcommand{\at}{\char64}                     % the @ character in running text
\newcommand{\contraction}[1]{\overbracket{#1}}           % contraction
% names of integral blocks
\newcommand{\oooo}{\text{oooo}}
\newcommand{\ooov}{\text{ooov}}
\newcommand{\oovo}{\text{oovo}}
\newcommand{\oovv}{\text{oovv}}
\newcommand{\ovov}{\text{ovov}}
\newcommand{\ovvo}{\text{ovvo}}
\newcommand{\ovvv}{\text{ovvv}}
\newcommand{\vvvv}{\text{vvvv}}
\newcommand{\vvvo}{\text{vvvo}}
\newcommand{\ovoo}{\text{ovoo}}
\newcommand{\vovv}{\text{vovv}}


% Running heads: book.cls upper-cases chapter and section titles, so long English titles overflow the page.
% Keep their case (otherwise the logic of book.cls); a chapter title that still does not fit can be given
% a short running-head title by defining \headtitle@<chapter number>.
% Page numbers always at the bottom centre (chapter openings too); running heads at the top without numbers.
\usepackage{fancyhdr}
\pagestyle{fancy}
\fancyhf{}
\fancyhead[LE]{\small\itshape\leftmark}
\fancyhead[RO]{\small\itshape\rightmark}
\fancyfoot[C]{\thepage}
\renewcommand{\headrulewidth}{0pt}
\fancypagestyle{plain}{\fancyhf{}\fancyfoot[C]{\thepage}\renewcommand{\headrulewidth}{0pt}}

\makeatletter
\renewcommand{\chaptermark}[1]{\markboth{\ifnum\c@secnumdepth>\m@ne\if@mainmatter\@chapapp\ \thechapter.\ \ \fi\fi
  \ifcsname headtitle@\thechapter\endcsname\csname headtitle@\thechapter\endcsname\else#1\fi}{}}
\renewcommand{\sectionmark}[1]{\markright{\ifnum\c@secnumdepth>\z@\thesection.\ \ \fi#1}}
\makeatother

\theoremstyle{definition}
\newtheorem{definition}{Definition}[chapter]
\newtheorem{theorem}{Theorem}[chapter]
\newtheorem{lemma}[theorem]{Lemma}
\newtheorem{remark}{Remark}[chapter]
\newtheorem{example}{Example}[chapter]

% ---------------- code listings and boxes ----------------
\lstdefinestyle{py}{
  language=Python, basicstyle=\ttfamily\small, keywordstyle=\color{blue!60!black}\bfseries,
  commentstyle=\color{green!40!black}, stringstyle=\color{red!50!black}, showstringspaces=false,
  columns=fullflexible, breaklines=true, frame=single, rulecolor=\color{gray!50},
  backgroundcolor=\color{gray!4}, numbers=left, numberstyle=\tiny\color{gray}, xleftmargin=2em,
  escapeinside={(*@}{@*)}
}
\lstset{style=py}

\newtcolorbox{codemap}[1][]{colback=blue!3,colframe=blue!40!black,title={Formula $\leftrightarrow$ code},fonttitle=\bfseries,#1}
\newtcolorbox{pitfall}[1][]{colback=red!3,colframe=red!50!black,title={Pitfall},fonttitle=\bfseries,#1}
\newtcolorbox{summary}[1][]{colback=green!3,colframe=green!40!black,title={Summary},fonttitle=\bfseries,#1}

\title{\Huge From Hartree--Fock to Coupled Cluster:\\[0.5em]\Large Theory Notes}
\author{}
\date{}
\hypersetup{pdftitle={From Hartree–Fock to Coupled Cluster: Theory Notes}}

\begin{document}
\frontmatter
\maketitle
\tableofcontents

\mainmatter
\chapter{The Variational Principle and Linear Variation}
\label{ch:variational}

All of electronic-structure theory rests on two things: the eigenvalue problem of the Schrödinger equation, and the variational principle used to approximate it when it cannot be solved exactly. HF, CI and MCSCF are variational methods; CC and perturbation theory are not, but their starting point (the HF reference) is. This chapter derives both from scratch and ends with the single line \code{scipy.linalg.eigh} in the code.

\textbf{Parallel reading.} “qc” pp.~30--32 (variational principle) and pp.~19--23 (basis expansion and matrix representation, change of basis); MEST pp.~2--5 (linear variation and the secular equation). Textbooks: Szabo--Ostlund §1.3 (the variational method; §1.3.2 the linear variational problem) and §1.1.5--1.1.7 (change of basis, the eigenvalue problem, functions of matrices---$S^{-1/2}$ is a function of a matrix), §3.4.5 (orthogonalization); Helgaker §4.2 (the variational principle).

\section{A Complete Set of Eigenfunctions and Expansions}

The non-relativistic electronic Hamiltonian in the Born--Oppenheimer approximation (atomic units) is
\begin{align}\label{eq:Hel}
\hat{H} &= -\sum_i \frac{1}{2} \nabla_i^2 - \sum_i \sum_A \frac{Z_A}{|r_i - R_A|} + \sum_{i<j} \frac{1}{|r_i - r_j|}
\end{align}
Let its eigenstates be $|\Phi_{\alpha}\rangle$, with eigenvalues in ascending order:
\begin{align*}
\hat{H} |\Phi_\alpha\rangle &= E_\alpha |\Phi_\alpha\rangle, \quad \alpha = 0, 1, 2, \ldots, \quad E_0 \le E_1 \le E_2 \le \ldots
\end{align*}
Because $\hat{H}$ is Hermitian, i.e.\ $\langle \Phi|\hat{H}|\Psi\rangle = \langle \Psi|\hat{H}|\Phi\rangle^*$ for any two states, its eigenvalues are real and eigenstates belonging to different eigenvalues are orthogonal. Degenerate eigenstates are not constrained in this way, but they can be Gram--Schmidt orthogonalized within the degenerate subspace. So one can always choose
\begin{align*}
\langle \Phi_\alpha | \Phi_\beta \rangle = \delta_{\alpha\beta}
\end{align*}
Whether this set of eigenstates is complete depends on the space we work in. Actual quantum-chemical calculations take place in the finite-dimensional space spanned by a finite basis, where $\hat{H}$ is a Hermitian matrix whose eigenvectors always form a complete orthonormal basis; this chapter is written for that finite-dimensional case. (For the true Hamiltonian in continuous space the bound states alone are not complete: the continuum must be added, and the sums below become sums plus integrals; the conclusions of the variational principle are unchanged.) The identity operator can then be written as
\begin{align}\label{eq:completeness}
\sum_{\alpha} |\Phi_\alpha\rangle\langle\Phi_\alpha| &= \hat{1}
\end{align}

Take any trial function $|\tilde{\Phi}\rangle$ satisfying the same boundary conditions and insert the identity to its left:
\begin{align*}
|\tilde{\Phi}\rangle &= \sum_{\alpha} |\Phi_\alpha\rangle \langle\Phi_\alpha|\tilde{\Phi}\rangle \\
&= \sum_{\alpha} |\Phi_\alpha\rangle C_\alpha, \quad C_\alpha = \langle\Phi_\alpha|\tilde{\Phi}\rangle
\end{align*}

\section{The Variational Principle}

\begin{theorem}[Variational principle]
For any normalized trial function $|\tilde{\Phi}\rangle$ satisfying the boundary conditions, $\langle\tilde{\Phi}|\hat{H}|\tilde{\Phi}\rangle \ge E_0$, with equality if and only if $|\tilde{\Phi}\rangle$ is the ground state (or lies in the degenerate ground-state subspace).
\end{theorem}

\begin{proof}
Idea: expand $|\tilde{\Phi}\rangle$ in the eigenstates; then $\langle\tilde{\Phi}|\hat{H}|\tilde{\Phi}\rangle$ is a weighted average of the eigenvalues and cannot be below the smallest one, $E_0$.

Step 1: express the normalization condition in terms of the expansion coefficients. Insert the completeness relation \eqref{eq:completeness} twice into $\langle\tilde{\Phi}|\tilde{\Phi}\rangle$:
\begin{align*}
1 &= \langle\tilde{\Phi}|\tilde{\Phi}\rangle \\
&= \sum_{\alpha \beta} \langle\tilde{\Phi}|\Phi_\alpha\rangle\langle\Phi_\alpha|\Phi_\beta\rangle\langle\Phi_\beta|\tilde{\Phi}\rangle \\
&= \sum_{\alpha \beta} \langle\tilde{\Phi}|\Phi_\alpha\rangle \delta_{\alpha \beta} \langle\Phi_\beta|\tilde{\Phi}\rangle \\
&= \sum_{\alpha} |\langle\Phi_\alpha|\tilde{\Phi}\rangle|^2 = \sum_{\alpha} |C_\alpha|^2
\end{align*}
The third equality uses orthonormality, $\langle \Phi_\alpha | \Phi_\beta \rangle = \delta_{\alpha\beta}$; the fourth sums over $\beta$, leaving only the term $\beta = \alpha$, and uses $\langle\tilde{\Phi}|\Phi_\alpha\rangle = \langle\Phi_\alpha|\tilde{\Phi}\rangle^*$.

Step 2: expand the energy expectation value in the same way, now with $\hat{H}$ in the middle. With $\hat{H}$ acting to the right on an eigenstate, the matrix element in the middle is
\begin{align*}
\langle \Phi_\alpha|\hat{H}|\Phi_\beta\rangle &= E_\beta \langle \Phi_\alpha|\Phi_\beta\rangle = E_\beta \delta_{\alpha\beta}
\end{align*}
i.e.\ $\hat{H}$ is diagonal in the basis of its own eigenstates. Hence
\begin{align*}
\langle\tilde{\Phi}|\hat{H}|\tilde{\Phi}\rangle &= \sum_{\alpha \beta} \langle\tilde{\Phi}|\Phi_\alpha\rangle\langle\Phi_\alpha|\hat{H}|\Phi_\beta\rangle\langle\Phi_\beta|\tilde{\Phi}\rangle \\
&= \sum_{\alpha \beta} \langle\tilde{\Phi}|\Phi_\alpha\rangle E_\beta \delta_{\alpha \beta} \langle\Phi_\beta|\tilde{\Phi}\rangle \\
&= \sum_{\alpha} E_\alpha |C_\alpha|^2
\end{align*}
Step 3: subtract $E_0$ from the result of Step 2, writing $E_0$ as $E_0 \sum_\alpha |C_\alpha|^2$ by Step 1:
\begin{align*}
\langle\tilde{\Phi}|\hat{H}|\tilde{\Phi}\rangle - E_0 &= \sum_\alpha (E_\alpha - E_0) |C_\alpha|^2 \ge 0
\end{align*}
Every term is a product of two non-negative numbers: $E_\alpha - E_0 \ge 0$ ($E_0$ is the lowest eigenvalue) and $|C_\alpha|^2 \ge 0$.

Step 4: equality. A sum of non-negative terms vanishes if and only if every term vanishes, i.e.\ for every $\alpha$ with $E_\alpha > E_0$ we have $C_\alpha = 0$. Then $|\tilde{\Phi}\rangle$ contains only components with eigenvalue $E_0$, $\hat{H} |\tilde{\Phi}\rangle = E_0 |\tilde{\Phi}\rangle$, and it is the ground state (or a vector in the degenerate ground-state subspace); conversely, the ground state obviously gives equality.
\end{proof}

For an unnormalized trial function (only $|\tilde{\Phi}\rangle \neq 0$ is required), use the Rayleigh quotient:
\begin{align}\label{eq:rayleigh}
E[\tilde{\Phi}] &= \frac{\langle\tilde{\Phi}|\hat{H}|\tilde{\Phi}\rangle}{\langle\tilde{\Phi}|\tilde{\Phi}\rangle} \ge E_0
\end{align}
with the same condition for equality as in the theorem above.

\begin{proof}[Proof of \eqref{eq:rayleigh}]
The denominator is positive: the expansion of Step 1 above holds for any $|\tilde{\Phi}\rangle$, $\langle\tilde{\Phi}|\tilde{\Phi}\rangle = \sum_\alpha |C_\alpha|^2$, and for $|\tilde{\Phi}\rangle \neq 0$ the coefficients do not all vanish, so $\langle\tilde{\Phi}|\tilde{\Phi}\rangle > 0$.

$E[\tilde{\Phi}]$ does not depend on the length or the phase of $|\tilde{\Phi}\rangle$: for any complex $\lambda \neq 0$ the coefficient in a bra is complex conjugated, $\langle\lambda\tilde{\Phi}| = \lambda^* \langle\tilde{\Phi}|$, so
\begin{align*}
E[\lambda\tilde{\Phi}] &= \frac{|\lambda|^2 \langle\tilde{\Phi}|\hat{H}|\tilde{\Phi}\rangle}{|\lambda|^2 \langle\tilde{\Phi}|\tilde{\Phi}\rangle} \\
&= E[\tilde{\Phi}]
\end{align*}
With $\lambda = \langle\tilde{\Phi}|\tilde{\Phi}\rangle^{-1/2}$, $|\Phi'\rangle = \lambda |\tilde{\Phi}\rangle$ is normalized ($\langle\Phi'|\Phi'\rangle = |\lambda|^2 \langle\tilde{\Phi}|\tilde{\Phi}\rangle = 1$), and the theorem applies to it:
\begin{align*}
E[\tilde{\Phi}] &= E[\Phi'] = \langle\Phi'|\hat{H}|\Phi'\rangle \ge E_0
\end{align*}
Equality holds if and only if $|\Phi'\rangle$ is a ground state, and $|\tilde{\Phi}\rangle$ differs from $|\Phi'\rangle$ only by a non-zero factor.
\end{proof}

\section{Linear Variation: Orthonormal Basis}

The most common trial function is a linear combination of a fixed set of functions $\{|\Psi_i\rangle\}$ ($i = 1, ..., N$), with the combination coefficients as variational parameters:
\begin{align*}
|\tilde{\Phi}\rangle &= \sum_{i=1}^N c_i |\Psi_i\rangle
\end{align*}
First assume that $\{|\Psi_i\rangle\}$ is orthonormal: $\langle\Psi_i|\Psi_j\rangle = \delta_{ij}$. Define the Hamiltonian matrix
\begin{align*}
H_{ij} &= \langle \Psi_i|\hat{H}|\Psi_j\rangle, \quad H_{ji}^* = \langle \Psi_j|\hat{H}|\Psi_i\rangle^* \\
&= \langle \Psi_i|\hat{H}^\dagger|\Psi_j\rangle = H_{ij}
\end{align*}
i.e.\ $H$ is a Hermitian matrix. This uses the definition of the adjoint: the adjoint (Hermitian conjugate) of an operator $\hat{A}$ is the operator $\hat{A}^\dagger$ satisfying, for all $|\Phi\rangle$ and $|\Psi\rangle$,
\begin{align*}
\langle\Phi|\hat{A}^\dagger|\Psi\rangle &= \langle\Psi|\hat{A}|\Phi\rangle^*
\end{align*}
and an operator with $\hat{A}^\dagger = \hat{A}$ is called Hermitian. Taking $|\Phi\rangle = |\Psi_i\rangle$ and $|\Psi\rangle = |\Psi_j\rangle$ gives the second equality in the line for $H_{ji}^*$ above, and $\hat{H}^\dagger = \hat{H}$ gives the last one. For matrices the adjoint is the conjugate transpose, $(A^\dagger)_{ij} = A_{ji}^*$, so $H_{ji}^* = H_{ij}$ means $H^\dagger = H$; a real Hermitian matrix is a real symmetric matrix.

In terms of the coefficients, the normalization condition and the energy are
\begin{align*}
\langle\tilde{\Phi}|\tilde{\Phi}\rangle &= \sum_{ij} c_i^* c_j \langle\Psi_i|\Psi_j\rangle = \sum_i |c_i|^2 = 1
\end{align*}
\begin{align*}
\langle\tilde{\Phi}|\hat{H}|\tilde{\Phi}\rangle = \sum_{ij} c_i^* c_j H_{ij}
\end{align*}

Now we look for the coefficients that minimize the energy. We cannot simply differentiate $\langle\tilde{\Phi}|\hat{H}|\tilde{\Phi}\rangle = \sum_{ij} c_i^* c_j H_{ij}$: it is the energy only when $\sum_i |c_i|^2 = 1$, and scaling all coefficients by $t$ multiplies it by $|t|^2$, so looking for its lowest point on its own only drives the coefficients to zero or to infinity. The remedy is the Rayleigh quotient of the previous section:
\begin{align*}
E(c) &= \frac{c^\dagger H c}{c^\dagger c}
\end{align*}
The denominator already contains the normalization, and $E(c)$ does not depend on the length of $c$ (proof of \eqref{eq:rayleigh}), so every coefficient can be varied freely. Take real coefficients first. Multiplying by the denominator, $E(c)\, c^{\mathsf T} c = c^{\mathsf T} H c$; differentiating with respect to $c_k$ (product rule):
\begin{align*}
\frac{\partial E}{\partial c_k} (c^{\mathsf T} c) + E \frac{\partial}{\partial c_k}(c^{\mathsf T} c) &= \frac{\partial}{\partial c_k}(c^{\mathsf T} H c)
\end{align*}
When differentiating, $c_k$ appears in two places of the double sum:
\begin{align*}
\frac{\partial}{\partial c_k} \left(\sum_{ij} c_i c_j H_{ij}\right) &= \sum_{ij} (\delta_{ik} c_j + c_i \delta_{jk}) H_{ij} \\
&= \sum_j c_j H_{kj} + \sum_i c_i H_{ik} = 2 \sum_j c_j H_{kj}
\end{align*}
The last step uses the real symmetry $H_{ik} = H_{ki}$; likewise $\partial (c^{\mathsf T} c) / \partial c_k = 2 c_k$. At the lowest point (in general, at a stationary point) $\frac{\partial E}{\partial c_k} = 0$, so for every $k$ we have $E c_k = \sum_j H_{kj} c_j$; together this is the matrix eigenvalue equation
\begin{align}\label{eq:Hc-Ec}
H c &= E c
\end{align}
For complex coefficients, differentiating with respect to the real and imaginary parts of each coefficient gives the same equation. Multiplying \eqref{eq:Hc-Ec} from the left by $c^\dagger$ and using $c^\dagger c = 1$ gives $E = c^\dagger H c$: the eigenvalue is the energy of this trial function.

MEST pp.~3--4 handle the normalization constraint with the method of Lagrange multipliers and obtain the same equation. That method becomes indispensable only for the HF equations (where the constraints are the orthonormality of all pairs of orbitals, which cannot be removed by a single quotient); it is introduced at the beginning of Chapter 3.

A Hermitian matrix $H$ has $N$ real eigenvalues and an orthonormal set of eigenvectors $c^{\alpha}$:
\begin{align*}
H c^{\alpha} &= E_{\alpha} c^{\alpha}, \quad (c^{\alpha})^{\dagger} c^{\beta} = \sum_i c_i^{\alpha *} c_i^{\beta} = \delta_{\alpha\beta}
\end{align*}
Arranging all eigenvectors into the matrix $C = [c^{0}, c^{1}, \ldots, c^{N-1}]$ and the eigenvalues into the diagonal matrix $E = \text{diag}(E_0, \ldots, E_{N-1})$,
\begin{align*}
H C & = C E, \quad C^\dagger C = 1
\end{align*}
The corresponding approximate states $|\tilde{\Phi}_\alpha\rangle = \sum_i c_i^\alpha |\Psi_i\rangle$ are also orthonormal and diagonalize $\hat{H}$ within this subspace:
\begin{align*}
\langle\tilde{\Phi}_\alpha|\tilde{\Phi}_\beta\rangle &= \sum_{ij} c_i^{\alpha *} c_j^\beta \delta_{ij} = \delta_{\alpha \beta}
\end{align*}
\begin{align*}
\langle\tilde{\Phi}_\alpha|\hat{H}|\tilde{\Phi}_\beta\rangle &= \sum_{ij} c_i^{\alpha *} H_{ij} c_j^\beta \\
&= (c^\alpha)^\dagger H c^\beta = E_\beta \delta_{\alpha \beta}
\end{align*}
The lowest eigenvalue (written $\tilde{E}_0$ below, to distinguish it from the exact $E_0$) is the best upper bound to the ground-state energy attainable in this function space.

\section{Linear Variation: Non-orthogonal Basis}

The basis functions of quantum chemistry (Gaussian atomic orbitals, non-orthogonal determinants) are generally not orthogonal. Let the overlap matrix be
\begin{align*}
S_{ij} &= \langle \Psi_i|\Psi_j\rangle, \quad S = S^\dagger, \quad S \ \text{positive definite}
\end{align*}
(Positive definiteness: for any non-zero $c$, $c^\dagger S c = \langle \tilde{\Phi}|\tilde{\Phi}\rangle > 0$, as long as the basis functions are linearly independent.)

As in the previous section, take the extremum of the Rayleigh quotient. With real coefficients,
\begin{align*}
E(c) &= \frac{c^{\mathsf T} H c}{c^{\mathsf T} S c} \Rightarrow E c^{\mathsf T} S c = c^{\mathsf T} H c
\end{align*}
Differentiate with respect to $c_k$ (both sides, with the product rule):
\begin{align*}
\frac{\partial E}{\partial c_k} (c^{\mathsf T} S c) + E \frac{\partial}{\partial c_k}(c^{\mathsf T} S c) &= \frac{\partial}{\partial c_k}(c^{\mathsf T} H c)
\end{align*}
The same calculation as in the previous section gives $\frac{\partial (c^{\mathsf T} S c)}{\partial c_k} = 2 (S c)_k$ and $\frac{\partial (c^{\mathsf T} H c)}{\partial c_k} = 2 (H c)_k$. At the extremum $\frac{\partial E}{\partial c_k} = 0$, so
\begin{align}\label{eq:HcESc}
2 E (S c)_k &= 2 (H c)_k \quad\Rightarrow\quad H c = E S c
\end{align}
This is the \textbf{generalized eigenvalue problem}. Arranging all solutions into a matrix:
\begin{align*}
H C &= S C E, \quad C^\dagger S C = 1
\end{align*}
The latter is the new orthonormality condition: the overlap between the approximate states $|\tilde{\Phi}_\alpha\rangle = \sum_i C_{i \alpha} |\Psi_i\rangle$ is exactly $(C^\dagger S C)_{\alpha \beta}$.

The Roothaan equations $F C = S C \epsilon$ (Chapter 4) have exactly this form, except that $F$ depends on $C$.

\section{Numerical Solution of the Generalized Eigenvalue Problem}

\subsection{Reduction to an Ordinary Eigenvalue Problem}

The standard way to solve $H c = E S c$: find a matrix $X$ such that
\begin{align}\label{eq:XSX}
X^\dagger S X &= 1
\end{align}
then set $c = X c'$, substitute, and multiply from the left by $X^\dagger$:
\begin{align*}
H X c' & = E S X c' \Rightarrow (X^\dagger H X) c' = E (X^\dagger S X) c' = E c'
\end{align*}
Defining $H' = X^\dagger H X$, the problem becomes the ordinary Hermitian eigenvalue problem $H' c' = E c'$. After solving for $c'$, back-substitute $c = X c'$. One can check that the orthonormality condition is satisfied automatically:
\begin{align*}
C^\dagger S C &= C'^\dagger X^\dagger S X C' = C'^\dagger C' = 1
\end{align*}

\subsection{Symmetric Orthogonalization (Löwdin)}

$S$ is a Hermitian positive-definite matrix and can be diagonalized: $S = U s U^\dagger$, $s = \mathrm{diag}(s_1, \ldots, s_N)$, $s_i > 0$, $U^\dagger U = 1$. Define
\begin{align*}
S^{-1/2} &= U s^{-1/2} U^\dagger, \quad s^{-1/2} = \text{diag}(s_1^{-1/2}, \ldots, s_N^{-1/2})
\end{align*}
Take $X = S^{-1/2}$ and check \eqref{eq:XSX}:
\begin{align*}
X^\dagger S X &= U s^{-1/2} U^\dagger U s U^\dagger U s^{-1/2} U^\dagger \\
&= U s^{-1/2} s s^{-1/2} U^\dagger \\
&= U U^\dagger \\
&= 1
\end{align*}
($S^{-1/2}$ is itself Hermitian, so $X^\dagger = X$.) This is exactly how MEST p.~14 turns $F C = S C \epsilon$ into $F' C' = C' \epsilon$:
\begin{align*}
F' &= S^{-1/2} F S^{-1/2}, \quad C = S^{-1/2} C'
\end{align*}

\subsection{Canonical Orthogonalization and Linear Dependence}

Another choice is $X = U s^{-1/2}$ (without the final $U^\dagger$):
\begin{align*}
X^\dagger S X &= s^{-1/2} U^\dagger U s U^\dagger U s^{-1/2} \\
&= s^{-1/2} s s^{-1/2} \\
&= 1
\end{align*}
Its advantage is that very small $s_i$ can be \emph{dropped}. When the basis contains very diffuse functions or atoms are close together, the basis functions become nearly linearly dependent; the smallest eigenvalue of $S$ can be as small as $10^{-8}$ or smaller, $s_i^{-1/2}$ then reaches $10^4$, and rounding errors are amplified by a factor of $10^4$ or more. Canonical orthogonalization keeps only the columns with $s_i > \tau$ (we take $\tau = 10^{-8}$):
\begin{align*}
X &= U_{:,\,\text{keep}}\, s_{\text{keep}}^{-1/2}, \quad X \ \text{is an}\ N \times M \ \text{matrix}, \quad M \le N
\end{align*}
Then $X^\dagger S X = 1_M$ still holds (only the dimension is now $M$), and the new eigenvalue problem is $M \times M$. The number of molecular orbitals $M$ can be smaller than the number of atomic orbitals $N$.

Implementation: \code{solve} in \code{hf\_to\_cc/scf.py} solves $F C = S C \epsilon$ with \code{scipy.linalg.eigh}.

\begin{summary}
\begin{itemize}
  \item Variational principle: the energy expectation value of any trial function is an upper bound to the ground-state energy.
  \item Linear variation $\Rightarrow$ $H c = E S c$; for an orthonormal basis $S = 1$.
  \item Numerically, an orthogonalizing matrix $X$ ($X^{\mathsf T} S X = 1$) reduces it to an ordinary eigenvalue problem; canonical orthogonalization can drop linearly dependent directions and is more robust.
\end{itemize}
\end{summary}

\chapter{Slater Determinants and the Slater--Condon Rules}
\label{ch:slater}

The HF energy expression, CI matrix elements, and the numerators and denominators of perturbation theory are all governed by one set of rules: the matrix elements of one- and two-body operators between two determinants. This chapter states them in first quantization (coordinate representation).

\textbf{Parallel reading.} “qc” pp.~55--58 (energy of a two-electron single determinant, expanded term by term) and “Adv qc” pp.~26--32 (for general $N$, the Slater--Condon rules derived one by one with permutation sums and arranged in the form of Szabo Tables 2.5 and 2.6); MEST pp.~5--10, 15. Textbooks: Szabo--Ostlund §2.2 (Slater determinants), §2.3.3--2.3.4 (general rules for matrix elements and their derivation); ESQC collection Ch.~24 (Almlöf et al., Slater determinants); CTDL Vol.~III, Ch.~XV §B-4 and §C-4--C-5 (reduced density operators, consequences of exchange) and Complement A$_{\rm XVI}$ §2 (the exchange hole of fermions). Chapter 7 rederives the single- and double-excitation cases in second quantization (“Adv qc” pp.~46--47, MEST p.~27).

\section{Antisymmetry and Slater Determinants}

One-electron wave functions (spin orbitals) are written $\chi_{i}(x)$, where $x = (r, \omega)$ comprises the spatial and the spin coordinates, and are taken to be orthonormal:
\begin{align*}
\langle \chi_i|\chi_j\rangle &= \int dx \, \chi_i^*(x) \chi_j(x) = \delta_{ij}
\end{align*}

\subsection{Two Electrons}

Subtract the exchanged term from the product $\chi_i(x_1) \chi_j(x_2)$:
\begin{align*}
\Phi(x_1, x_2) &= \frac{1}{\sqrt{2}} [\chi_i(x_1) \chi_j(x_2) - \chi_j(x_1) \chi_i(x_2)] \\
&= \frac{1}{\sqrt{2}} \det\begin{pmatrix} \chi_i(x_1) & \chi_j(x_1) \\ \chi_i(x_2) & \chi_j(x_2) \end{pmatrix}
\end{align*}
(Read the last matrix as a 2×2 determinant: rows are electrons, columns are orbitals.) Antisymmetry: exchanging $x_1, x_2$ swaps two rows of the determinant, which changes its sign. If $i = j$ the determinant has two equal columns and vanishes---this is the Pauli exclusion principle. Below, $\chi_i(1)$ is shorthand for $\chi_i(x_1)$.

\subsection{N Electrons}

In general a Slater determinant is written
\begin{align*}
\Phi(x_1, \dots, x_N) &= \frac{1}{\sqrt{N!}} \det[\chi_{k_j}(x_i)] = |\chi_{k_1} \chi_{k_2} \dots \chi_{k_N}\rangle
\end{align*}

\section{Energy of a Two-Electron Determinant}

Before the general $N$-electron results, the energy of the simplest determinant, two electrons in the spin orbitals $\chi_i$ and $\chi_j$, is worked out term by term (MEST pp.~5--6, “qc” pp.~55--57). For two electrons the Hamiltonian \eqref{eq:Hel} is a one-electron operator for each electron (kinetic energy plus nuclear attraction) plus the repulsion between the two:
\begin{align*}
\hat{H} &= \left[ -\frac{1}{2} \nabla_1^2 - \sum_A \frac{Z_A}{|r_1 - R_A|} \right] + \left[ -\frac{1}{2} \nabla_2^2 - \sum_A \frac{Z_A}{|r_2 - R_A|} \right] + \frac{1}{|r_1 - r_2|} \\
&= h(1) + h(2) + r_{12}^{-1}
\end{align*}
$h(1)$ acts only on $x_1$, and $h(2)$ only on $x_2$. The energy of $\Phi(x_1, x_2) = \frac{1}{\sqrt{2}} [\chi_i(1) \chi_j(2) - \chi_j(1) \chi_i(2)]$ splits into three parts:
\begin{align*}
E &= \langle \Phi|\hat{H}|\Phi\rangle = \langle \Phi|h(1)|\Phi\rangle + \langle \Phi|h(2)|\Phi\rangle + \langle \Phi|r_{12}^{-1}|\Phi\rangle
\end{align*}

\subsection{One-Electron Part}

\begin{align*}
\langle \Phi|h(1)|\Phi\rangle &= \int dx_1 dx_2 \left[ \frac{1}{\sqrt{2}} (\chi_i(1) \chi_j(2) - \chi_j(1) \chi_i(2)) \right]^* h(1) \left[ \frac{1}{\sqrt{2}} (\chi_i(1) \chi_j(2) - \chi_j(1) \chi_i(2)) \right] \\
&= \frac{1}{2} \int dx_1 dx_2\, [\chi_i(1) \chi_j(2) - \chi_j(1) \chi_i(2)]^* h(1) [\chi_i(1) \chi_j(2) - \chi_j(1) \chi_i(2)] \\
&= \frac{1}{2} \int dx_1 dx_2\, \chi_i^*(1) \chi_j^*(2) h(1) \chi_i(1) \chi_j(2) \eqnote{multiply out} \\
&\quad - \frac{1}{2} \int dx_1 dx_2\, \chi_i^*(1) \chi_j^*(2) h(1) \chi_j(1) \chi_i(2) \\
&\quad - \frac{1}{2} \int dx_1 dx_2\, \chi_j^*(1) \chi_i^*(2) h(1) \chi_i(1) \chi_j(2) \\
&\quad + \frac{1}{2} \int dx_1 dx_2\, \chi_j^*(1) \chi_i^*(2) h(1) \chi_j(1) \chi_i(2) \\
&= \frac{1}{2} \int dx_2\, \chi_j^*(2) \chi_j(2) \int dx_1\, \chi_i^*(1) h(1) \chi_i(1) \eqnote{$h(1)$ acts on $x_1$ only} \\
&\quad - \frac{1}{2} \int dx_2\, \chi_j^*(2) \chi_i(2) \int dx_1\, \chi_i^*(1) h(1) \chi_j(1) \\
&\quad - \frac{1}{2} \int dx_2\, \chi_i^*(2) \chi_j(2) \int dx_1\, \chi_j^*(1) h(1) \chi_i(1) \\
&\quad + \frac{1}{2} \int dx_2\, \chi_i^*(2) \chi_i(2) \int dx_1\, \chi_j^*(1) h(1) \chi_j(1) \\
&= \frac{1}{2} [\langle i|h|i\rangle + \langle j|h|j\rangle] \eqnote{orthonormality: the $x_2$ integrals are $1, 0, 0, 1$}
\end{align*}
where $\langle p|h|q\rangle = \int dx_1\, \chi_p^*(1) h(1) \chi_q(1)$. The same steps for $h(2)$, whose $x_1$ integrals now give the overlaps, yield
\begin{align*}
\langle \Phi|h(2)|\Phi\rangle &= \frac{1}{2} [\langle j|h|j\rangle + \langle i|h|i\rangle]
\end{align*}
so that
\begin{align*}
\langle \Phi|h(1) + h(2)|\Phi\rangle &= \langle i|h|i\rangle + \langle j|h|j\rangle
\end{align*}

\subsection{Two-Electron Part}

\begin{align*}
\langle \Phi|r_{12}^{-1}|\Phi\rangle &= \int dx_1 dx_2 \left[ \frac{1}{\sqrt{2}} (\chi_i(1) \chi_j(2) - \chi_j(1) \chi_i(2)) \right]^* r_{12}^{-1} \left[ \frac{1}{\sqrt{2}} (\chi_i(1) \chi_j(2) - \chi_j(1) \chi_i(2)) \right] \\
&= \frac{1}{2} \int dx_1 dx_2\, \chi_i^*(1) \chi_j^*(2) r_{12}^{-1} \chi_i(1) \chi_j(2) \\
&\quad - \frac{1}{2} \int dx_1 dx_2\, \chi_i^*(1) \chi_j^*(2) r_{12}^{-1} \chi_j(1) \chi_i(2) \\
&\quad - \frac{1}{2} \int dx_1 dx_2\, \chi_j^*(1) \chi_i^*(2) r_{12}^{-1} \chi_i(1) \chi_j(2) \\
&\quad + \frac{1}{2} \int dx_1 dx_2\, \chi_j^*(1) \chi_i^*(2) r_{12}^{-1} \chi_j(1) \chi_i(2) \\
&= \frac{1}{2} \int dx_1 dx_2\, \chi_i^*(1) \chi_j^*(2) r_{12}^{-1} \chi_i(1) \chi_j(2) \\
&\quad - \frac{1}{2} \int dx_1 dx_2\, \chi_i^*(1) \chi_j^*(2) r_{12}^{-1} \chi_j(1) \chi_i(2) \\
&\quad - \frac{1}{2} \int dx_1 dx_2\, \chi_j^*(2) \chi_i^*(1) r_{12}^{-1} \chi_i(2) \chi_j(1) \eqnote{$x_1 \leftrightarrow x_2$, $r_{21} = r_{12}$; $=$ 2nd term} \\
&\quad + \frac{1}{2} \int dx_1 dx_2\, \chi_j^*(2) \chi_i^*(1) r_{12}^{-1} \chi_j(2) \chi_i(1) \eqnote{$x_1 \leftrightarrow x_2$; $=$ 1st term} \\
&= \frac{1}{2} \left[ 2 \int dx_1 dx_2\, \chi_i^*(1) \chi_j^*(2) r_{12}^{-1} \chi_i(1) \chi_j(2) - 2 \int dx_1 dx_2\, \chi_i^*(1) \chi_j^*(2) r_{12}^{-1} \chi_j(1) \chi_i(2) \right] \\
&= \int dx_1 dx_2\, \chi_i^*(1) \chi_i(1) r_{12}^{-1} \chi_j^*(2) \chi_j(2) - \int dx_1 dx_2\, \chi_i^*(1) \chi_j(1) r_{12}^{-1} \chi_j^*(2) \chi_i(2) \\
&= \phys{ij}{ij} - \phys{ij}{ji} = \asym{ij}{ij}
\end{align*}
where
\begin{align*}
\phys{pq}{rs} &= \int dx_1 dx_2\, \chi_p^*(1) \chi_q^*(2) r_{12}^{-1} \chi_r(1) \chi_s(2) \\
\asym{pq}{rs} &= \phys{pq}{rs} - \phys{pq}{sr}
\end{align*}
The Coulomb term $\phys{ij}{ij} = \int dx_1 dx_2\, |\chi_i(1)|^2 r_{12}^{-1} |\chi_j(2)|^2$ is the classical repulsion between the charge clouds $|\chi_i|^2$ and $|\chi_j|^2$; it is also the only two-electron term that the plain product $\chi_i(1) \chi_j(2)$ (the Hartree product) gives. The exchange term $\phys{ij}{ji}$ has no classical counterpart: it comes from the antisymmetry of the determinant.

Altogether, the energy of the two-electron determinant is
\begin{align*}
E &= \langle i|h|i\rangle + \langle j|h|j\rangle + \asym{ij}{ij}
\end{align*}
With $N$ electrons the same expansion gives one such one-electron term for each occupied spin orbital and one Coulomb-minus-exchange term for each pair of occupied spin orbitals; the next two sections state these results.


\section{Matrix Elements of One-Electron Operators}

Let $O_1 = \sum_{i=1}^N h(i)$, where $h$ acts only on the coordinates of electron $i$ (e.g.\ kinetic energy plus nuclear attraction, the $h_{\text{core}}$ of Chapter 4). Its expectation value in the determinant $\Phi$ is the sum of the expectation values of $h$ in the occupied spin orbitals:
\begin{align}\label{eq:sc-1e-diag}
\langle \Phi|O_1|\Phi\rangle &= \sum_{m \in \text{occ}} \langle m|h|m\rangle
\end{align}

\section{Matrix Elements of Two-Electron Operators}

Let $O_2 = \sum_{i<j} v(i, j)$ with $v(i, j) = \frac{1}{r_{ij}}$. Each pair of occupied spin orbitals $m, n$ contributes a Coulomb integral minus an exchange integral:
\begin{align*}
\phys{mn}{mn} &= \int dx_1 dx_2 \chi_m^*(1) \chi_n^*(2) r_{12}^{-1} \chi_m(1) \chi_n(2)
\end{align*}
\begin{align*}
\phys{mn}{nm} &= \int dx_1 dx_2 \chi_m^*(1) \chi_n^*(2) r_{12}^{-1} \chi_n(1) \chi_m(2)
\end{align*}
Summed over all pairs,
\begin{align}\label{eq:sc-2e-diag}
\langle \Phi|O_2|\Phi\rangle &= \sum_{m<n} [\phys{mn}{mn} - \phys{mn}{nm}] = \sum_{m<n} \asym{mn}{mn} \\
&= \frac{1}{2} \sum_{mn} \asym{mn}{mn} \notag
\end{align}
The last step: $\asym{mn}{mn}$ is symmetric in $m, n$ ($\asym{nm}{nm} = \asym{mn}{mn}$, by exchanging the two integration variables), and for $m = n$, $\asym{mm}{mm} = \phys{mm}{mm} - \phys{mm}{mm} = 0$, so $\sum_{m<n}$ equals $\frac{1}{2} \sum_{mn}$. This is very important: it frees us from worrying about ``self-interaction''---the $m = n$ self-interaction in the Coulomb term is cancelled exactly by the exchange term.

Combining the one- and two-body parts, the energy of a single determinant is
\begin{align}
E[\Phi] &= \sum_m \langle m|h|m\rangle + \frac{1}{2} \sum_{mn} \asym{mn}{mn} \notag \\
&= \sum_m h_{mm} + \frac{1}{2} \sum_{mn} [\phys{mn}{mn} - \phys{mn}{nm}] \label{eq:E-det}
\end{align}
This is the HF energy expression; Chapter 3 makes it variational.

\section{The Slater--Condon Rules}

CI needs matrix elements between different determinants. Let $|K\rangle$ and $|L\rangle$ be arranged in \emph{maximum coincidence} (equal orbitals in equal positions; the reordering brings a sign, discussed below).

\textbf{Case 1: $|K\rangle = |L\rangle$.} Already obtained above:
\begin{align*}
\langle K|O_1|K\rangle &= \sum_{m \in K} \langle m|h|m\rangle, \\
\langle K|O_2|K\rangle &= \frac{1}{2} \sum_{m,n \in K} \asym{mn}{mn}
\end{align*}

\textbf{Case 2: one orbital different.} $|K\rangle = |\dots m \dots\rangle$, $|L\rangle = |\dots p \dots\rangle$ ($m$ replaced by $p$, everything else equal):
\begin{align*}
\langle K|\hat{O}_1|L\rangle &= \langle m|h|p\rangle, \quad \langle K|\hat{O}_2|L\rangle = \sum_{n \in K} \asym{mn}{pn}
\end{align*}

\textbf{Case 3: two orbitals different.} $m, n \rightarrow p, q$:
\begin{align*}
\langle K|\hat{O}_1|L\rangle &= 0, \quad \langle K|\hat{O}_2|L\rangle = \asym{mn}{pq}
\end{align*}
A one-body operator can ``repair'' only one differing orbital, and the other differing orbital gives a zero overlap; a two-body operator repairs exactly two.

\textbf{Case 4: three or more orbitals different}: all matrix elements vanish.

\begin{pitfall}
\textbf{Signs.} The formulas above assume that the two determinants are already in maximum coincidence. If the original orbital order of $|L\rangle$ is different, its columns must first be swapped into place, with a factor $-1$ for every swap.
\end{pitfall}

\section{Spin Integration: From Spin Orbitals to Spatial Orbitals}

Restricted spin orbitals are spatial orbitals times spin functions:
\begin{align*}
\chi_{2i-1}(x) &= \psi_i(r) \alpha(\omega), \quad \chi_{2i}(x) = \psi_i(r) \beta(\omega) \\
\langle \alpha|\alpha\rangle &= \langle \beta|\beta\rangle = 1, \langle \alpha|\beta\rangle = 0
\end{align*}
For the integral $\int dx = \int dr \sum_{\omega}$ the spin part can be integrated out first. For a one-electron integral with $\chi_p = \psi_p \sigma_p$, $\chi_q = \psi_q \sigma_q$:
\begin{align*}
\langle \chi_p|h|\chi_q\rangle &= (\psi_p|h|\psi_q) \langle \sigma_p|\sigma_q\rangle = h_{pq} \delta_{\sigma_p \sigma_q}
\end{align*}
A two-electron integral (chemists' notation, the two orbitals of the same electron together):
\begin{align*}
[\chi_p \chi_q|\chi_r \chi_s] &= (pq|rs) \delta_{\sigma_p \sigma_q} \delta_{\sigma_r \sigma_s}
\end{align*}
That is: the two spin orbitals of the same electron must have the same spin. In physicists' notation $\phys{pq}{rs} = [pr|qs]$:
\begin{align*}
\langle \chi_p \chi_q|\chi_r \chi_s\rangle &= (pr|qs) \delta_{\sigma_p \sigma_r} \delta_{\sigma_q \sigma_s}
\end{align*}

\subsection{Energy of a Closed-Shell Determinant}

Closed shell: $\frac{N}{2}$ spatial orbitals $\psi_i$, each occupied by one $\alpha$ and one $\beta$ electron. Split the spin-orbital sums of \eqref{eq:E-det} into sums over spatial orbitals and spins. One-body part:
\begin{align*}
\sum_m \langle m|h|m\rangle &= \sum_{i}^{N/2} [h_{ii} + h_{ii}] = 2 \sum_i h_{ii}
\end{align*}
Two-body Coulomb part: $\phys{mn}{mn} = (mm|nn) \delta_{\sigma_m \sigma_m} \delta_{\sigma_n \sigma_n} = (ii|jj)$ is non-zero for all 4 spin combinations,
\begin{align*}
\frac{1}{2} \sum_{mn} \phys{mn}{mn} &= \frac{1}{2} \cdot 4 \sum_{ij} (ii|jj) = 2 \sum_{ij} (ii|jj)
\end{align*}
Exchange part: $\phys{mn}{nm} = (mn|nm) \delta_{\sigma_m \sigma_n} \delta_{\sigma_n \sigma_m}$ is non-zero only for the two combinations $\alpha \alpha$ and $\beta \beta$,
\begin{align*}
\frac{1}{2} \sum_{mn} \phys{mn}{nm} &= \frac{1}{2} \cdot 2 \sum_{ij} (ij|ji) = \sum_{ij} (ij|ji)
\end{align*}
So the closed-shell energy is
\begin{align}\label{eq:E-rhf-mo}
E_{\text{RHF}} &= 2 \sum_i h_{ii} + \sum_{ij} [2 (ii|jj) - (ij|ji)] = 2 \sum_i h_{ii} + \sum_{ij} (2 J_{ij} - K_{ij})
\end{align}
$J_{ij} = (ii|jj)$ is the Coulomb integral and $K_{ij} = (ij|ji)$ the exchange integral. This agrees exactly with MEST p.~15; Chapter 4 rewrites it in terms of the AO density matrix.

\begin{summary}
\begin{itemize}
\item Single-determinant energy $E = \sum_m h_{mm} + \frac{1}{2} \sum_{mn} \asym{mn}{mn}$; the $m = n$ self-interaction cancels automatically.
\item Slater--Condon: determinants differing by 0, 1 or 2 orbitals give three formulas, more differences give zero; mind the reordering sign.
\item Spin integration: the two orbitals of the same electron must have the same spin; closed shell $E = 2 \sum h_{ii} + \sum (2J - K)$.
\end{itemize}
\end{summary}

\chapter{The Hartree--Fock Equations}
\label{ch:hf}

The previous chapter gave the energy of a single determinant, \eqref{eq:E-det}. The HF method is: among all single determinants, find the one with the lowest energy. This chapter derives the HF equations by functional variation, following the order of MEST pp.~7--12 and writing out every step; it then discusses canonicalization, orbital energies and Brillouin's theorem.

\textbf{Parallel reading.} “qc” pp.~59--61 derive the HF equations with Lagrange multipliers along the same route as this chapter; p.~66 discusses orbital energies; “Adv qc” p.~31 and MEST p.~27 cover Brillouin's theorem. Textbooks: Szabo--Ostlund §3.1--3.3 (§3.2.2 minimization of the single-determinant energy, §3.3.2 Brillouin's theorem); Helgaker §10.1--10.5 (derivation with exponentially parametrized orbital rotations); ESQC collection Ch.~25. ESQC exercise 77.6 derives the HF conditions from the first-order variation of a determinant and makes a good exercise.

\section{Statement of the Problem}

The energy is a functional of the $N$ spin orbitals:
\begin{align*}
E[\{\chi\}] &= \sum_i \langle i|h|i\rangle + \frac{1}{2} \sum_{ij} \asym{ij}{ij}
\end{align*}
The constraint is orthonormality of the orbitals: $\langle i | j \rangle = \delta_{ij}$. Since $\langle j|i\rangle = \langle i|j\rangle^*$, the independent real conditions are $N$ normalization conditions plus $N(N-1)/2$ complex orthogonality conditions (each with a real and an imaginary part), $N^2$ in total; for real orbitals there are $N(N+1)/2$. Writing one multiplier for every pair $(i, j)$ does no harm; the relation between the extra multipliers is given in the next section (the multiplier matrix is Hermitian).

The standard way to find an extremum under constraints is the method of Lagrange multipliers: give each constraint an undetermined coefficient (a multiplier), subtract ``multiplier $\times$ constraint'' from the energy to obtain a new function $L$, and then find its stationary points by varying all variables freely, as if there were no constraints; the values of the multipliers are fixed at the end by the constraints themselves. The reason: at the lowest point under the constraints, no change that respects the constraints alters the energy (to first order), so the energy can change only along directions that violate the constraints, and subtracting ``multiplier $\times$ constraint'' exactly cancels that change. The linear variation of Chapter 1 has a single constraint (normalization), which the Rayleigh quotient removes; the multiplier method gives the same equation (real coefficients):
\begin{align*}
L &= c^{\mathsf T} H c - E (c^{\mathsf T} c - 1) \\
\frac{\partial L}{\partial c_k} &= 2 (H c)_k - 2 E c_k = 0 \quad \Rightarrow \quad H c = E c
\end{align*}
The multiplier $E$ is the energy. Here there are $N^2$ constraints and no such quotient is available, so this section handles them directly with multipliers (the alternative is to write the orbitals as unitary rotations of reference orbitals, which satisfy the constraints automatically). The HF multipliers also have a clear meaning: as shown below, they are the matrix elements of the Fock operator between occupied orbitals, and after diagonalization they are the orbital energies. Introduce multipliers $\epsilon_{ji}$:
\begin{align}\label{eq:hf-lagrangian}
L[\{\chi\}] &= E[\{\chi\}] - \sum_{ij} \epsilon_{ji} (\langle i|j\rangle - \delta_{ij})
\end{align}
and require $L$ to be stationary under arbitrary variations of the orbitals: $\delta L = 0$.

\section{The Lagrange Multipliers Form a Hermitian Matrix}

$E$ is real, and $L$ must be real as well (otherwise ``minimization'' makes no sense). Since $E$ and $\delta_{ij}$ are real, $\sum_{ij} \epsilon_{ji} \langle i|j\rangle$ must be real:
\begin{align*}
\sum_{ij} \epsilon_{ji} \langle i|j\rangle &= \left( \sum_{ij} \epsilon_{ji} \langle i|j\rangle \right)^* \\
&= \sum_{ij} \epsilon_{ji}^* \langle j|i\rangle
\end{align*}
Swap the dummy indices $i \leftrightarrow j$ on the right:
\begin{align*}
\sum_{ij} \epsilon_{ji}^* \langle j|i\rangle &= \sum_{ij} \epsilon_{ij}^* \langle i|j\rangle
\end{align*}
So $\sum_{ij} (\epsilon_{ji} - \epsilon_{ij}^*) \langle i|j\rangle = 0$. This holds for arbitrary orbitals (arbitrary values of $\langle i | j \rangle$), hence
\begin{align*}
\epsilon_{ji} &= \epsilon_{ij}^*, \quad \text{i.e.} \quad \epsilon^\dagger = \epsilon
\end{align*}
MEST pp.~7--8 argue slightly differently (requiring the pair of terms $\epsilon_{ij} \langle i|j\rangle + \epsilon_{ji} \langle j|i\rangle$ to be real), with the same conclusion.

\section{First-Order Variation of the Energy}

Change one orbital $\chi_m$ into $\chi_m + \delta \chi_m$ and keep the others fixed (the last subsection of this section shows that varying all orbitals at once gives the same condition). Keep only first-order terms.

\subsection{One-Body Part}

\begin{align*}
\delta \sum_i \langle i|h|i\rangle &= \langle m + \delta m|h|m + \delta m\rangle - \langle m|h|m\rangle \\
&= \langle \delta m|h|m\rangle + \langle m|h|\delta m\rangle + O(\delta^2)
\end{align*}
By the Hermiticity of $h$, $\langle m|h|\delta m\rangle = \langle h m|\delta m\rangle = \langle \delta m|h|m\rangle^*$, so
\begin{align*}
\delta \left( \sum_i \langle i|h|i\rangle \right) &= \langle \delta m|h|m\rangle + \text{c.c.}
\end{align*}
(c.c.\ denotes the complex conjugate of the preceding terms.)

\subsection{Two-Body Part}

This is where mistakes are easiest, so write it term by term. In the two-body energy $\frac{1}{2} \sum_{ij} \asym{ij}{ij}$ the orbital $m$ appears in three classes of terms: $i = m, j \neq m$; $i \neq m, j = m$; $i = j = m$. The last class, $\asym{mm}{mm} = 0$, vanishes identically and stays zero after the variation (the antisymmetrized integral satisfies this for any $\chi_m$), so it can be ignored. In the notation of MEST pp.~9--10, split into three parts:
\begin{align*}
\frac{1}{2} \sum_{ij} \asym{ij}{ij} &= \frac{1}{2} \sum_{i \neq m} \sum_{j \neq m} \asym{ij}{ij} \\
&\quad + \frac{1}{2} \sum_{j \neq m} \asym{mj}{mj} + \frac{1}{2} \sum_{i \neq m} \asym{im}{im}
\end{align*}
The first part does not depend on $\chi_m$. The third part equals the second after exchanging the integration variables of the two electrons ($x_1 \leftrightarrow x_2$):
\begin{align*}
\asym{im}{im} &= \int dx_1 dx_2 \chi_i^*(1) \chi_m^*(2) r_{12}^{-1} [\chi_i(1) \chi_m(2) - \chi_m(1) \chi_i(2)] \\
& = \int dx_2 dx_1 \chi_i^*(2) \chi_m^*(1) r_{12}^{-1} [\chi_i(2) \chi_m(1) - \chi_m(2) \chi_i(1)] = \asym{mi}{mi}
\end{align*}
So the part that depends on $m$ is $\sum_{j \neq m} \asym{mj}{mj}$. Now vary. In $\asym{mj}{mj}$, $\chi_m$ appears twice (the $\chi_m^*$ on the left and the $\chi_m$ on the right), and the first-order variation has two terms:
\begin{align*}
\delta \asym{mj}{mj} &= \asym{\delta m j}{m j} + \asym{m j}{\delta m j}
\end{align*}
The second term is the complex conjugate of the first:
\begin{align*}
\asym{m j}{\delta m j} &= \int dx_1 dx_2 \chi_m^*(1) \chi_j^*(2) r_{12}^{-1} [\delta \chi_m(1) \chi_j(2) - \chi_j(1) \delta \chi_m(2)] \\
&= \left( \int dx_1 dx_2 \delta \chi_m^*(1) \chi_j^*(2) r_{12}^{-1} [\chi_m(1) \chi_j(2) - \chi_j(1) \chi_m(2)] \right)^* \\
&= \asym{\delta m j}{m j}^*
\end{align*}
(Second equality: conjugate the bracket as a whole and exchange $x_1, x_2$ in the exchange term, then compare with the first term.) Hence
\begin{align*}
\delta \left( \frac{1}{2} \sum_{ij} \asym{ij}{ij} \right) &= \sum_j \asym{\delta m j}{m j} + \text{c.c.}
\end{align*}
where the $j = m$ term has been added back, since $\asym{\delta m m}{m m} = \phys{\delta m m}{m m} - \phys{\delta m m}{m m} = 0$.

\subsection{Constraint Part}

\begin{align*}
\delta \sum_{ij} \epsilon_{ji} \langle i|j\rangle &= \sum_j \epsilon_{jm} \langle \delta m|j\rangle + \sum_i \epsilon_{mi} \langle i|\delta m\rangle \\
&= \sum_j \epsilon_{jm} \langle \delta m|j\rangle + \text{c.c.}
\end{align*}
The second equality uses $\epsilon_{mi} \langle i|\delta m\rangle = (\epsilon_{mi}^* \langle\delta m|i\rangle)^* = (\epsilon_{im} \langle\delta m|i\rangle)^*$ and renames the dummy index $i$ to $j$.

\subsection{Varying All Orbitals at Once}

MEST pp.~7--8 vary all orbitals simultaneously, $\chi_i \rightarrow \chi_i + \delta \chi_i$ for $i = 1, \dots, N$, and arrive at the same condition. To first order,
\begin{align*}
\delta L &= \sum_i \delta \langle i|h|i\rangle + \frac{1}{2} \sum_{ij} \delta [\phys{ij}{ij} - \phys{ij}{ji}] - \sum_{ij} \epsilon_{ji}\, \delta \langle i|j\rangle \eqnote{$\delta_{ij}$ is a constant}
\end{align*}
The one-body part is the one-body result above, summed over $i$: $\sum_i \delta \langle i|h|i\rangle = \sum_i \langle \delta i|h|i\rangle + \sum_i \langle i|h|\delta i\rangle = \sum_i \langle \delta i|h|i\rangle + \text{c.c.}$ In the two-body part both orbitals of every integral vary, which gives eight first-order terms:
\begin{align*}
\frac{1}{2} \sum_{ij} \delta [\phys{ij}{ij} - \phys{ij}{ji}] &= \frac{1}{2} \sum_{ij} \phys{\delta i\, j}{ij} + \frac{1}{2} \sum_{ij} \phys{i\, \delta j}{ij} + \frac{1}{2} \sum_{ij} \phys{ij}{\delta i\, j} + \frac{1}{2} \sum_{ij} \phys{ij}{i\, \delta j} \\
&\quad - \frac{1}{2} \sum_{ij} \phys{\delta i\, j}{ji} - \frac{1}{2} \sum_{ij} \phys{i\, \delta j}{ji} - \frac{1}{2} \sum_{ij} \phys{ij}{\delta j\, i} - \frac{1}{2} \sum_{ij} \phys{ij}{j\, \delta i}
\end{align*}
The second term equals the first, and the sixth equals the fifth (rename the summation indices $i \leftrightarrow j$, then the integration variables $x_1 \leftrightarrow x_2$):
\begin{align*}
\sum_{ij} \phys{i\, \delta j}{ij} &= \sum_{ij} \phys{j\, \delta i}{ji} = \sum_{ij} \int dx_1 dx_2\, \chi_j^*(1) \delta \chi_i^*(2) r_{12}^{-1} \chi_j(1) \chi_i(2) \eqnote{$i \leftrightarrow j$} \\
&= \sum_{ij} \int dx_1 dx_2\, \delta \chi_i^*(1) \chi_j^*(2) r_{12}^{-1} \chi_i(1) \chi_j(2) = \sum_{ij} \phys{\delta i\, j}{ij} \eqnote{$x_1 \leftrightarrow x_2$}
\end{align*}
\begin{align*}
\sum_{ij} \phys{i\, \delta j}{ji} &= \sum_{ij} \phys{j\, \delta i}{ij} = \sum_{ij} \int dx_1 dx_2\, \chi_j^*(1) \delta \chi_i^*(2) r_{12}^{-1} \chi_i(1) \chi_j(2) \eqnote{$i \leftrightarrow j$} \\
&= \sum_{ij} \int dx_1 dx_2\, \delta \chi_i^*(1) \chi_j^*(2) r_{12}^{-1} \chi_j(1) \chi_i(2) = \sum_{ij} \phys{\delta i\, j}{ji} \eqnote{$x_1 \leftrightarrow x_2$}
\end{align*}
The terms with $\delta$ in the ket are complex conjugates of those with $\delta$ in the bra: $\phys{ij}{\delta i\, j} = \phys{\delta i\, j}{ij}^*$, $\phys{ij}{i\, \delta j} = \phys{i\, \delta j}{ij}^*$, and, after renaming $i \leftrightarrow j$, $\phys{ji}{\delta i\, j} = \phys{\delta i\, j}{ji}^*$, $\phys{ji}{i\, \delta j} = \phys{i\, \delta j}{ji}^*$. So the eight terms collapse to two and their conjugates:
\begin{align*}
\frac{1}{2} \sum_{ij} \delta [\phys{ij}{ij} - \phys{ij}{ji}] &= \frac{1}{2} \sum_{ij} [2 \phys{\delta i\, j}{ij} - 2 \phys{\delta i\, j}{ji}] + \text{c.c.} \\
&= \sum_{ij} \asym{\delta i\, j}{ij} + \text{c.c.}
\end{align*}
The constraint part, with the second-order term $\langle \delta i|\delta j\rangle$ dropped:
\begin{align*}
\sum_{ij} \epsilon_{ji}\, \delta \langle i|j\rangle &= \sum_{ij} \epsilon_{ji} \langle \delta i|j\rangle + \sum_{ij} \epsilon_{ji} \langle i|\delta j\rangle = \sum_{ij} \epsilon_{ji} \langle \delta i|j\rangle + \sum_{ij} \epsilon_{ij} \langle j|\delta i\rangle \eqnote{$i \leftrightarrow j$} \\
&= \sum_{ij} \epsilon_{ji} \langle \delta i|j\rangle + \sum_{ij} \epsilon_{ji}^* \langle \delta i|j\rangle^* = \sum_{ij} \epsilon_{ji} \langle \delta i|j\rangle + \text{c.c.} \eqnote{$\epsilon_{ij} = \epsilon_{ji}^*$}
\end{align*}
Altogether,
\begin{align*}
\delta L &= \sum_i \Big[ \langle \delta i|h|i\rangle + \sum_j \asym{\delta i\, j}{ij} - \sum_j \epsilon_{ji} \langle \delta i|j\rangle \Big] + \text{c.c.} = 0
\end{align*}
Term $i$ of this sum is what the three parts above give for $m = i$. The $\delta \chi_i$ are independent, so setting all of them to zero except one, $\delta \chi_m$, leaves exactly the one-orbital variation: varying all orbitals at once gives the same condition, for every $m$, as varying them one at a time. The next section works it out.

\section{The Hartree--Fock Equations (Non-canonical Form)}

Combining the three parts:
\begin{align*}
\delta L &= \langle \delta m|h|m\rangle + \sum_j \asym{\delta m j}{m j} \\
&\quad - \sum_j \epsilon_{jm} \langle \delta m|j\rangle + \text{c.c.} = 0
\end{align*}
In integral form (with electron 1 as the ``varied electron''):
\begin{align*}
\delta L &= \int dx_1 \delta \chi_m^*(1) \left[ h(1) \chi_m(1) + \sum_j \int dx_2 \chi_j^*(2) r_{12}^{-1} \chi_j(2) \chi_m(1) \right. \\
& \left. - \sum_j \int dx_2 \chi_j^*(2) r_{12}^{-1} \chi_m(2) \chi_j(1) - \sum_j \epsilon_{jm} \chi_j(1) \right] + \text{c.c.} = 0
\end{align*}
Call the part before c.c.\ $Z$; then $\delta L = Z + Z^* = 2 \operatorname{Re} Z$. Replacing $\delta \chi_m$ by $i\, \delta\chi_m$ turns $Z$ into $-iZ$ and gives $\delta L = 2 \operatorname{Im} Z$; both vanish, so $Z = 0$ for every $\delta \chi_m$, and the bracket must vanish identically. Define the Coulomb and exchange operators:
\begin{align*}
\hat{J}_j(1) \chi(1) &= \left[ \int dx_2 \chi_j^*(2) r_{12}^{-1} \chi_j(2) \right] \chi(1)
\end{align*}
\begin{align*}
\hat{K}_j(1) \chi(1) &= \left[ \int dx_2 \chi_j^*(2) r_{12}^{-1} \chi(2) \right] \chi_j(1)
\end{align*}
$\hat{J}_j$ is local: it simply multiplies by the electrostatic potential of the charge distribution of electron $j$. $\hat{K}_j$ is non-local: it ``moves'' the function $\chi$ it acts on into the integral and then multiplies by $\chi_j$. Define the Fock operator
\begin{align}\label{eq:fock-op}
\hat{f}(1) &= h(1) + \sum_j [\hat{J}_j(1) - \hat{K}_j(1)]
\end{align}
The condition that the bracket vanishes becomes
\begin{align}\label{eq:hf-noncanonical}
\hat{f} |\chi_m\rangle &= \sum_j \epsilon_{jm} |\chi_j\rangle, \quad m = 1, \dots, N
\end{align}
These are the non-canonical HF equations. Multiplying from the left by $\langle \chi_i|$ shows $\epsilon_{im} = \langle i|\hat{f}|m\rangle$: the matrix of Lagrange multipliers is the matrix of the Fock operator between occupied orbitals.

\section{Unitary Invariance and Canonicalization}

\subsection{The Determinant Is Invariant Under Unitary Transformations of the Occupied Orbitals}

Apply a unitary transformation to the occupied orbitals, $\chi'_i = \sum_j \chi_j U_{ji}$, $U^\dagger U = 1$. The determinant is the determinant of the matrix $A_{ki} = \chi_i(x_k)$; after the transformation $A' = A U$ and
\begin{align*}
\det(A') &= \det(A) \det(U), \quad |\det U|^2 = \det(U^\dagger) \det(U) = \det(U^\dagger U) \\
&= 1 \quad \Rightarrow \quad \det U = e^{i \theta}
\end{align*}
So the new and old determinants differ only by a phase, and all physical quantities (energy, density) are unchanged.

MEST pp.~11--12 check this intuitively with a two-orbital example; here it is expanded in full. Take $U = \frac{1}{\sqrt{2}} \begin{pmatrix} 1 & 1 \\ -1 & 1 \end{pmatrix}$ ($\det U = \frac{1}{2} (1 + 1) = 1$); with $\chi'_i = \sum_j \chi_j U_{ji}$,
\begin{align*}
\chi'_1 &= \frac{1}{\sqrt{2}} (\chi_1 - \chi_2), \quad \chi'_2 = \frac{1}{\sqrt{2}} (\chi_1 + \chi_2)
\end{align*}
The new determinant (omitting the normalization factor $\frac{1}{\sqrt{2}}$) is
\begin{align*}
\chi'_1(1) \chi'_2(2) - \chi'_2(1) \chi'_1(2) &= \frac{1}{2} [(\chi_1(1) - \chi_2(1))(\chi_1(2) + \chi_2(2)) - (\chi_1(1) + \chi_2(1))(\chi_1(2) - \chi_2(2))]
\end{align*}
Expanding each product into four terms:
\begin{align*}
&= \frac{1}{2} [\chi_1(1) \chi_1(2) + \chi_1(1) \chi_2(2) - \chi_2(1) \chi_1(2) - \chi_2(1) \chi_2(2) \\
&\quad - \chi_1(1) \chi_1(2) + \chi_1(1) \chi_2(2) - \chi_2(1) \chi_1(2) + \chi_2(1) \chi_2(2)] \\
&= \chi_1(1) \chi_2(2) - \chi_2(1) \chi_1(2)
\end{align*}
The ``two electrons in the same orbital'' terms $\chi_1(1) \chi_1(2)$ and $\chi_2(1) \chi_2(2)$ cancel, and what remains is exactly the original determinant, consistent with $\det U = 1$. The shapes of the individual orbitals have changed completely, yet the determinant has not changed at all: the HF wave function is determined by the occupied space, not by a particular set of orbitals.

\subsection{The Fock Operator Is Invariant}

$\hat{J}$ and $\hat{K}$ depend on the orbitals only through $\sum_j |\chi_j\rangle\langle\chi_j|$. For the Coulomb operator:
\begin{align*}
\sum_j \hat{J}'_j(1) &= \sum_j \int d x_2 \, \chi'^*_j(2) r_{12}^{-1} \chi'_j(2) \\
&= \sum_j \sum_{kl} U_{kj}^* U_{lj} \int d x_2 \, \chi_k^*(2) r_{12}^{-1} \chi_l(2) \\
&= \sum_{kl} (U U^\dagger)_{lk} \int d x_2 \, \chi_k^*(2) r_{12}^{-1} \chi_l(2) \\
&= \sum_k \hat{J}_k(1)
\end{align*}
The derivation for the exchange operator is the same, except that the function $\chi$ being acted on enters the integral and $\chi'_j$ stays outside:
\begin{align*}
\sum_j \hat{K}'_j(1) \chi(1) &= \sum_j \left[ \int dx_2 \chi'^*_j(2) r_{12}^{-1} \chi(2) \right] \chi'_j(1) \\
&= \sum_{kl} (U U^\dagger)_{lk} \left[ \int dx_2 \chi_k^*(2) r_{12}^{-1} \chi(2) \right] \chi_l(1) \\
&= \sum_k \hat{K}_k(1) \chi(1)
\end{align*}
So $\hat{f}' = \hat{f}$.

\subsection{Transformation and Diagonalization of the Multiplier Matrix}

\begin{align*}
\epsilon'_{ij} &= \langle \chi'_i | \hat{f} | \chi'_j \rangle \\
&= \sum_{kl} U_{ki}^* \langle k | \hat{f} | l \rangle U_{lj} = (U^\dagger \epsilon U)_{ij}
\end{align*}
$\epsilon$ is Hermitian, so $U$ can always be chosen to make $U^\dagger \epsilon U$ diagonal. In this set of orbitals \eqref{eq:hf-noncanonical} becomes an eigenvalue equation:
\begin{align}\label{eq:hf-canonical}
\hat{f} | \chi_i \rangle &= \epsilon_i | \chi_i \rangle
\end{align}
These are the canonical HF equations, and $\epsilon_i$ are the orbital energies. The virtual orbitals can be treated the same way: by \eqref{eq:hf-noncanonical}, $\hat{f}$ maps every occupied orbital to a combination of occupied orbitals, i.e.\ the occupied space $O$ is invariant under $\hat{f}$; hence $\langle a|\hat{f}|i\rangle = 0$ for any $|a\rangle$ in the orthogonal complement of $O$, and since $\hat{f}$ is Hermitian also $\langle i|\hat{f}|a\rangle = 0$, so $\hat{f}$ maps the complement into itself as well. In a finite basis $\hat{f}$ therefore splits into uncoupled occupied and virtual blocks, which can be diagonalized separately, giving $\hat{f} |a\rangle = \epsilon_a |a\rangle$. Rotations within the virtual block do not change the reference determinant. This is where the $\epsilon_a$ in the MP2 denominators come from.

\section{Orbital Energies and the Total Energy}

Using \eqref{eq:fock-op}:
\begin{align*}
\epsilon_i &= \langle i | \hat{f} | i \rangle = \langle i | h | i \rangle + \sum_j [\phys{ij}{ij} - \phys{ij}{ji}] \\
&= h_{ii} + \sum_j \asym{ij}{ij}
\end{align*}
Orbital energies of virtual orbitals are defined in the same way, $\epsilon_a = \langle a|\hat{f}|a\rangle = h_{aa} + \sum_j \asym{aj}{aj}$: a virtual orbital ``feels'' the field of all $N$ occupied electrons. In terms of the occupied orbital energies, the HF energy is
\begin{align}\label{eq:ehf-eps}
E_{\text{HF}} &= \sum_i \epsilon_i - \frac{1}{2} \sum_{ij} \asym{ij}{ij} = \frac{1}{2} \sum_i (h_{ii} + \epsilon_i)
\end{align}
The total energy is \emph{not} the sum of orbital energies: the sum of orbital energies counts the electron--electron repulsion twice. The expression $E = \tr(h D) + \frac{1}{2} \tr(V_{\text{eff}} D)$ of Chapter 4 is exactly the AO version of \eqref{eq:ehf-eps}.

\section{Brillouin's Theorem}

The matrix element between a singly excited determinant $|\Phi_i^a\rangle$ ($\chi_i \rightarrow \chi_a$) and the HF determinant follows from Slater--Condon case 2:
\begin{align*}
\langle \Phi_i^a | \hat{H} | \Phi_0 \rangle &= \langle a | h | i \rangle + \sum_j \asym{aj}{ij} \\
&= \langle a | \hat{f} | i \rangle = f_{ai}
\end{align*}
At the HF solution the matrix elements of $\hat{f}$ between occupied and virtual orbitals vanish (obvious in the canonical case; in the non-canonical case \eqref{eq:hf-noncanonical} shows that $\hat{f}|i\rangle$ contains only occupied orbitals), so
\begin{align*}
\langle \Phi_i^a | \hat{H} | \Phi_0 \rangle &= 0 \quad (\text{Brillouin theorem})
\end{align*}

\section{The Closed-Shell RHF Equations}

For a closed shell, each of the $\frac{N}{2}$ spatial orbitals $\psi_i$ holds two electrons. Split the sum over spin orbitals $j$ in \eqref{eq:fock-op} into spatial orbitals and spins. Let the Fock operator act on an $\alpha$ spin orbital $\chi(x_1) = \psi(r_1) \alpha(\omega_1)$, with the spin orbital $\chi_j$ running over $\psi_j \alpha$ and $\psi_j \beta$. In the integral of the Coulomb operator the spin part of $\chi_j^*(2) \chi_j(2)$ is $\langle \sigma_j|\sigma_j\rangle = 1$, so both spins contribute the same spatial integral:
\begin{align*}
\hat{J}_{j \alpha} \chi(1) + \hat{J}_{j \beta} \chi(1) &= 2 \left[ \int dr_2 \psi_j^*(r_2) r_{12}^{-1} \psi_j(r_2) \right] \chi(1)
\end{align*}
In the integral of the exchange operator we have $\chi_j^*(2) \chi(2)$, whose spin part is $\langle \sigma_j|\alpha\rangle$, so only the $\sigma_j = \alpha$ term survives:
\begin{align*}
\hat{K}_{j \alpha} \chi(1) + \hat{K}_{j \beta} \chi(1) &= \left[ \int dr_2 \psi_j^*(r_2) r_{12}^{-1} \psi(r_2) \right] \psi_j(r_1) \alpha(\omega_1) + 0
\end{align*}
The result is the same for a $\beta$ spin $\chi$ (then the $\sigma_j = \beta$ term survives). The brackets on the right-hand sides are the $\hat{J}_j$ and $\hat{K}_j$ containing only spatial integrals, so
\begin{align}\label{eq:fock-rhf}
\hat{f}(r_1) &= h(r_1) + \sum_{j=1}^{N/2} [2 \hat{J}_j(r_1) - \hat{K}_j(r_1)]
\end{align}
where $\hat{J}_j$ and $\hat{K}_j$ now contain only spatial integrals. The spatial orbitals satisfy $\hat{f} \psi_i = \epsilon_i \psi_i$.

\begin{summary}
\begin{itemize}
\item Varying $E[\{\chi\}]$ under the orthonormality constraint $\Rightarrow$ $\hat{f} \chi_m = \sum_j \epsilon_{jm} \chi_j$, with $\epsilon$ Hermitian.
\item The determinant and $\hat{f}$ are invariant under unitary transformations of the occupied orbitals, so $\epsilon$ can be diagonalized $\Rightarrow$ the canonical HF equations.
\item $E_{\text{HF}} = \frac{1}{2} \sum_i (h_{ii} + \epsilon_i)$.
\item Brillouin's theorem: at the HF solution $\langle \Phi_i^a | \hat{H} | \Phi_0 \rangle = f_{ai} = 0$.
\end{itemize}
\end{summary}

\chapter{The Roothaan--Hall Equations and the RHF Program}
\label{ch:roothaan}

The HF equations of Chapter 3 are integro-differential equations. Expanding the orbitals in a finite basis of atomic orbitals turns them into matrix equations---the Roothaan--Hall equations. This chapter gives the Fock matrix, density matrix and energy in the AO representation.

\textbf{Parallel reading.} “qc” pp.~62--65: J and K matrix elements in the AO basis, the density matrix, the Roothaan equations, the SCF iteration and Löwdin orthogonalization, in one-to-one correspondence with this chapter; MEST pp.~13--16. Textbooks: Szabo--Ostlund §3.4 (§3.4.2 the Roothaan equations, §3.4.4 the Fock matrix, §3.4.5 orthogonalization of the basis, §3.4.6 the SCF procedure); Helgaker §10.6.1; ESQC collection Ch.~26. ESQC exercises 77.5 (the density matrix in LCAO) and 77.7 (derivative of the energy with respect to LCAO coefficients) make good exercises.

\section{Basis-Set Expansion}

Take $K$ atomic orbitals (basis functions) $\{\phi_\mu\}$ and write the spatial molecular orbitals as
\begin{align*}
\psi_i = \sum_{\mu=1}^K C_{\mu i} \phi_\mu
\end{align*}
Substitute into the closed-shell HF equation $\hat{f} \psi_i = \epsilon_i \psi_i$:
\begin{align*}
\hat{f} \sum_{\nu} C_{\nu i} \phi_\nu &= \epsilon_i \sum_{\nu} C_{\nu i} \phi_\nu
\end{align*}
Multiply from the left by $\phi_\mu^*$ and integrate:
\begin{align*}
\sum_{\nu} \langle \phi_\mu|\hat{f}|\phi_\nu\rangle C_{\nu i} &= \epsilon_i \sum_{\nu} \langle \phi_\mu|\phi_\nu\rangle C_{\nu i}
\end{align*}
Defining the Fock matrix $F_{\mu \nu} = \langle \phi_\mu|\hat{f}|\phi_\nu\rangle$ and the overlap matrix $S_{\mu \nu} = \langle \phi_\mu|\phi_\nu\rangle$, this reads
\begin{align*}
\sum_{\nu} F_{\mu \nu} C_{\nu i} &= \epsilon_i \sum_{\nu} S_{\mu \nu} C_{\nu i} \quad \Rightarrow \quad (F C)_{\mu i} = (S C \epsilon)_{\mu i}
\end{align*}
Last step: let $\epsilon$ be the diagonal matrix $\epsilon_{ji} = \epsilon_i \delta_{ji}$; then $\epsilon_i (S C)_{\mu i} = \sum_j (S C)_{\mu j} \epsilon_{ji} = (S C \epsilon)_{\mu i}$. Hence
\begin{align}\label{eq:roothaan}
F C = S C \epsilon
\end{align}
This is the generalized eigenvalue problem of Chapter 1, except that $F$ depends on $C$, so it has to be solved iteratively (SCF).

\section{The Fock Matrix in the AO Representation}

\subsection{One-Electron Part}

\begin{align*}
H^{\text{core}}_{\mu \nu} &= \langle \phi_\mu|h|\phi_\nu\rangle = T_{\mu \nu} + V_{\mu \nu}, \\
T_{\mu \nu} &= \langle \phi_\mu| -\frac{1}{2} \nabla^2 |\phi_\nu\rangle, \\
V_{\mu \nu} &= \langle \phi_\mu| -\sum_A \frac{Z_A}{|\mathbf{r} - \mathbf{R}_A|} |\phi_\nu\rangle
\end{align*}
In PySCF these are \code{mol.intor("int1e\_kin")} and \code{mol.intor("int1e\_nuc")}.

\subsection{Coulomb Part}

\begin{align*}
&J_{\mu \nu} = \sum_j^{N/2} \langle \phi_\mu|\hat{J}_j|\phi_\nu\rangle \\
&= \sum_j^{N/2} \int d\mathbf{r}_1 d\mathbf{r}_2 \phi_\mu^*(1) \phi_\nu(1) r_{12}^{-1} \psi_j^*(2) \psi_j(2)
\end{align*}
Expand $\psi_j$: $\psi_j^*(2) \psi_j(2) = \sum_{\lambda \sigma} C_{\lambda j}^* C_{\sigma j} \phi_\lambda^*(2) \phi_\sigma(2)$. Then
\begin{align*}
J_{\mu \nu} = \sum_{\lambda \sigma} \left[ \sum_j^{N/2} C_{\lambda j}^* C_{\sigma j} \right] \chem{\mu \nu}{\lambda \sigma}
\end{align*}
where $(\mu \nu|\lambda \sigma) = \int d\mathbf{r}_1 d\mathbf{r}_2 \phi_\mu^*(1) \phi_\nu(1) r_{12}^{-1} \phi_\lambda^*(2) \phi_\sigma(2)$ (chemists' notation). Define the density matrix
\begin{align}\label{eq:density}
D_{\sigma \lambda} = 2 \sum_j^{N/2} C_{\sigma j} C_{\lambda j}^* = 2 (C_{\text{occ}} C_{\text{occ}}^\dagger)_{\sigma \lambda}
\end{align}
(The factor 2 comes from the two electrons per spatial orbital. Note the index order: MEST p.~13 deliberately writes $P_{\sigma \lambda}$, opposite to the $\lambda \sigma$ of the integral, so that it is also correct for complex orbitals.) Then
\begin{align*}
J_{\mu \nu} = \frac{1}{2} \sum_{\lambda \sigma} \chem{\mu \nu}{\lambda \sigma} D_{\sigma \lambda}
\end{align*}
For real orbitals $D$ is symmetric, $D_{\sigma \lambda} = D_{\lambda \sigma}$. The closed-shell Fock operator contains $2 \hat{J}_j$, so the Coulomb term of the Fock matrix is
\begin{align*}
2 J_{\mu \nu} = \sum_{\lambda \sigma} \chem{\mu \nu}{\lambda \sigma} D_{\lambda \sigma}
\end{align*}

\subsection{Exchange Part}

\begin{align*}
&K_{\mu \nu} = \sum_j^{N/2} \langle \phi_\mu|\hat{K}_j|\phi_\nu\rangle \\
&= \sum_j^{N/2} \int d\mathbf{r}_1 d\mathbf{r}_2 \phi_\mu^*(1) \psi_j(1) r_{12}^{-1} \psi_j^*(2) \phi_\nu(2)
\end{align*}
(By the definition of $\hat{K}_j$: the function acted on, $\phi_\nu$, enters the integral and moves to electron 2, while $\psi_j$ stays with electron 1.) Expanding $\psi_j$ in the same way:
\begin{align*}
K_{\mu \nu} &= \sum_{\lambda \sigma} \left[ \sum_j C_{\lambda j} C_{\sigma j}^* \right] \chem{\mu \lambda}{\sigma \nu} \\
&= \frac{1}{2} \sum_{\lambda \sigma} \chem{\mu \lambda}{\sigma \nu} D_{\lambda \sigma}
\end{align*}
Altogether, the closed-shell Fock matrix is
\begin{align}\label{eq:fock-ao}
F_{\mu \nu} &= H^{\text{core}}_{\mu \nu} + \sum_{\lambda \sigma} D_{\lambda \sigma} \left[ \chem{\mu \nu}{\lambda \sigma} - \frac{1}{2} \chem{\mu \lambda}{\sigma \nu} \right]
\end{align}
This is identical to Szabo Eq.~(3.154) (Szabo's $P$ is our $D$).

In the program we define the Coulomb and exchange matrices (without factors) as
\begin{align*}
J[D]_{\mu \nu} &= \sum_{\lambda \sigma} \chem{\mu \nu}{\lambda \sigma} D_{\lambda \sigma}, \\
K[D]_{\mu \nu} &= \sum_{\lambda \sigma} \chem{\mu \lambda}{\sigma \nu} D_{\lambda \sigma}
\end{align*}
so that $F = H^{\text{core}} + J[D] - \frac{1}{2} K[D]$, which is exactly the line \code{F = h + J - 0.5 * K} of \code{rhf} in \code{hf\_to\_cc/scf.py}.

\section{Properties of the Density Matrix}

Let $C$ satisfy $C^{\mathsf T} S C = 1$ (real orbitals) and $D = 2 C_{\text{occ}} C_{\text{occ}}^{\mathsf T}$. The single-determinant energy depends only on $D$, not on the particular choice of occupied orbitals---this is the AO-representation form of the unitary invariance of Chapter 3: $D$ is unchanged under $C_{\text{occ}} \rightarrow C_{\text{occ}} U$.

\section{The Energy in the AO Representation}

The closed-shell energy of Chapter 2,
\begin{align*}
E_{\text{RHF}} &= 2 \sum_i h_{ii} + \sum_{ij} \left[ 2 \chem{ii}{jj} - \chem{ij}{ji} \right]
\end{align*}
is rewritten in the AO basis term by term (MEST pp.~15--16). With $\psi_i = \sum_\mu C_{\mu i} \phi_\mu$, each MO integral is a sum of AO integrals,
\begin{align*}
h_{ii} &= \sum_{\mu \nu} C_{\mu i}^* C_{\nu i} H^{\text{core}}_{\mu \nu} \\
\chem{ij}{kl} &= \sum_{\mu \nu \lambda \sigma} C_{\mu i}^* C_{\nu j} C_{\lambda k}^* C_{\sigma l} \chem{\mu \nu}{\lambda \sigma}
\end{align*}
and each sum over the occupied orbitals is half a density-matrix element, $\sum_i C_{\nu i} C_{\mu i}^* = \frac{1}{2} D_{\nu \mu}$ by \eqref{eq:density}.

\textbf{One-electron term.}
\begin{align*}
2 \sum_i h_{ii} &= 2 \sum_{\mu \nu} \sum_i C_{\mu i}^* C_{\nu i} H^{\text{core}}_{\mu \nu} \\
&= \sum_{\mu \nu} H^{\text{core}}_{\mu \nu} \Big( 2 \sum_i C_{\nu i} (C_{\text{occ}}^\dagger)_{i \mu} \Big) \eqnote{$C^\dagger_{i \mu} = C_{\mu i}^*$} \\
&= \sum_{\mu \nu} H^{\text{core}}_{\mu \nu}\, 2 (C_{\text{occ}} C_{\text{occ}}^\dagger)_{\nu \mu} \\
&= \sum_{\mu \nu} H^{\text{core}}_{\mu \nu} D_{\nu \mu} = \sum_\mu (H^{\text{core}} D)_{\mu \mu} = \tr(H^{\text{core}} D) \quad \text{by \eqref{eq:density}}
\end{align*}

\textbf{Coulomb term.}
\begin{align*}
2 \sum_{ij} \chem{ii}{jj} &= 2 \sum_{\mu \nu \lambda \sigma} \sum_{ij} C_{\mu i}^* C_{\nu i} C_{\lambda j}^* C_{\sigma j} \chem{\mu \nu}{\lambda \sigma} \\
&= 2 \sum_{\mu \nu \lambda \sigma} \chem{\mu \nu}{\lambda \sigma} \sum_i C_{\mu i}^* C_{\nu i} \sum_j C_{\lambda j}^* C_{\sigma j} \\
&= 2 \sum_{\mu \nu \lambda \sigma} \chem{\mu \nu}{\lambda \sigma} \left( \frac{1}{2} D_{\nu \mu} \right) \left( \frac{1}{2} D_{\sigma \lambda} \right) \\
&= \frac{1}{2} \sum_{\mu \nu \lambda \sigma} \chem{\mu \nu}{\lambda \sigma} D_{\nu \mu} D_{\sigma \lambda} \\
&= \frac{1}{2} \sum_{\mu \nu \lambda \sigma} \chem{\mu \nu}{\lambda \sigma} D_{\mu \nu} D_{\lambda \sigma} \eqnote{real orbitals: $D$ symmetric} \\
&= \frac{1}{2} \sum_{\mu \nu} J[D]_{\mu \nu} D_{\mu \nu} \eqnote{$J[D]_{\mu \nu} = \sum_{\lambda \sigma} \chem{\mu \nu}{\lambda \sigma} D_{\lambda \sigma}$} \\
&= \frac{1}{2} \sum_{\mu \nu} J[D]_{\mu \nu} D_{\nu \mu} = \frac{1}{2} \tr(J[D] D)
\end{align*}

\textbf{Exchange term.}
\begin{align*}
-\sum_{ij} \chem{ij}{ji} &= -\sum_{\mu \nu \lambda \sigma} \sum_{ij} C_{\mu i}^* C_{\nu j} C_{\lambda j}^* C_{\sigma i} \chem{\mu \nu}{\lambda \sigma} \\
&= -\sum_{\mu \nu \lambda \sigma} \chem{\mu \nu}{\lambda \sigma} \sum_i C_{\mu i}^* C_{\sigma i} \sum_j C_{\lambda j}^* C_{\nu j} \\
&= -\sum_{\mu \nu \lambda \sigma} \chem{\mu \nu}{\lambda \sigma} \left( \frac{1}{2} D_{\sigma \mu} \right) \left( \frac{1}{2} D_{\nu \lambda} \right) \\
&= -\frac{1}{4} \sum_{\mu \nu \lambda \sigma} \chem{\mu \nu}{\lambda \sigma} D_{\sigma \mu} D_{\nu \lambda} \\
&= -\frac{1}{4} \sum_{\mu \nu \lambda \sigma} \chem{\mu \nu}{\lambda \sigma} D_{\mu \sigma} D_{\nu \lambda} \eqnote{real orbitals} \\
&= -\frac{1}{4} \sum_{\mu \sigma} K[D]_{\mu \sigma} D_{\mu \sigma} \eqnote{$K[D]_{\mu \sigma} = \sum_{\nu \lambda} \chem{\mu \nu}{\lambda \sigma} D_{\nu \lambda}$} \\
&= -\frac{1}{4} \sum_{\mu \sigma} K[D]_{\mu \sigma} D_{\sigma \mu} = -\frac{1}{4} \tr(K[D] D)
\end{align*}
MEST p.~16 then writes $K[D]$ with the same summed pair $\lambda \sigma$ on $D$ as $J[D]$:
\begin{align*}
K[D]_{\mu \sigma} &= \sum_{\nu \lambda} \chem{\mu \nu}{\lambda \sigma} D_{\nu \lambda} = \sum_{\nu \lambda} \chem{\mu \nu}{\sigma \lambda} D_{\nu \lambda} \eqnote{real orbitals: $\chem{\mu \nu}{\lambda \sigma} = \chem{\mu \nu}{\sigma \lambda}$}
\end{align*}
and, renaming the summed $\nu, \lambda$ to $\lambda, \sigma$ and the free $\sigma$ to $\nu$,
\begin{align*}
K[D]_{\mu \nu} &= \sum_{\lambda \sigma} \chem{\mu \lambda}{\nu \sigma} D_{\lambda \sigma} \\
J[D]_{\mu \nu} &= \sum_{\lambda \sigma} \chem{\mu \nu}{\lambda \sigma} D_{\lambda \sigma}
\end{align*}
With the AO indices written $i, j, k, l$ (unrelated to the occupied orbitals above), each piece is one \code{np.einsum} call (MEST pp.~15--16):
\begin{center}
\begin{tabular}{ll}
\toprule
formula & \code{np.einsum} \\
\midrule
$\tr(H^{\text{core}} D) = \sum_{ij} H^{\text{core}}_{ij} D_{ji}$ & \code{np.einsum("ij,ji->", hcore, dm)} \\
$\tr(J[D] D) = \sum_{ij} J[D]_{ij} D_{ji}$ & \code{np.einsum("ij,ji->", j, dm)} \\
$\tr(K[D] D) = \sum_{ij} K[D]_{ij} D_{ji}$ & \code{np.einsum("ij,ji->", k, dm)} \\
$J[D]_{ij} = \sum_{kl} \chem{ij}{kl} D_{kl}$ & \code{np.einsum("ijkl,kl->ij", eri, dm)} \\
$K[D]_{ij} = \sum_{kl} \chem{ik}{jl} D_{kl}$ & \code{np.einsum("ikjl,kl->ij", eri, dm)} \\
\bottomrule
\end{tabular}
\end{center}
The program \code{hf\_to\_cc/scf.py} (box below) uses the same two contractions, with the letters \code{"mnls"} for \code{"ijkl"} and \code{"mlns"} for \code{"ikjl"}, and evaluates the energy in the combined form \mbox{$\frac{1}{2} \tr[(H + F) D]$}.

Adding the three terms and the nuclear repulsion $E_{nn}$, with $H = H^{\text{core}}$, $J = J[D]$ and $K = K[D]$:
\begin{align}\label{eq:e-ao}
E_{\text{RHF}} &= \tr(H D) + \frac{1}{2} \tr(J D) - \frac{1}{4} \tr(K D) + E_{nn} \\
&= \tr(H D) + \frac{1}{2} \tr(V_{\text{eff}} D) + E_{nn}, \quad V_{\text{eff}} = J - \frac{1}{2} K
\end{align}
Equivalently $E = \frac{1}{2} \tr[(H + F) D] + E_{nn}$, the AO version of \eqref{eq:ehf-eps}.

Implementation: \code{rhf} and \code{scf\_diis} (with the DIIS of Chapter 5) in \code{hf\_to\_cc/scf.py}.

\begin{codemap}
\begin{lstlisting}[frame=none, numbers=none, xleftmargin=0pt, backgroundcolor={}]
J = np.einsum("mnls, ls->mn", eri, d_o)
K = np.einsum("mlns, ls->mn", eri, d_o)
F = h + J - 0.5 * K
E_hf = 0.5 * np.einsum("mn, mn->", d_o, h+F) + E_nuc
\end{lstlisting}
\code{eri[m, n, l, s]} is $\chem{\mu\nu}{\lambda\sigma}$ and \code{d\_o} is $D$. The first two lines are $J[D]$ and $K[D]$ (for real orbitals $\chem{\mu\lambda}{\nu\sigma} = \chem{\mu\lambda}{\sigma\nu}$), the third is the Fock matrix \eqref{eq:fock-ao}, the last is $E = \frac{1}{2} \tr[(H + F) D] + E_{nn}$.
\end{codemap}

\begin{summary}
\begin{itemize}
\item The basis expansion turns the HF equations into $F C = S C \epsilon$, with $F = H + J[D] - \frac{1}{2} K[D]$ and $D = 2 C_{\text{occ}} C_{\text{occ}}^{\mathsf T}$.
\item $E = \tr(HD) + \frac{1}{2} \tr(JD) - \frac{1}{4} \tr(KD) + E_{nn}$.
\end{itemize}
\end{summary}

\chapter{SCF Convergence and DIIS}
\label{ch:scf}

With the Roothaan equations alone and no convergence acceleration, methanol and benzene simply cannot be computed. This chapter starts from ``SCF is a fixed-point iteration'', analyses why it oscillates, and then presents DIIS.

\textbf{Parallel reading.} ESQC collection Ch.~27 (Almlöf, Ahlrichs, Klopper): §27.1 level shifting and damping, §27.2 DIIS, §27.4 guessing initial orbitals. The experience reported there: DIIS needs only the last 4--6 steps; use level shifting for the first few steps and then switch to DIIS. Helgaker §10.6.2 interprets DIIS as a weighted average of the effective potentials of the previous steps (weights summing to 1, equivalent to averaging the density matrices) and notes that it is most effective when the naive SCF iteration oscillates between two potentials; §10.8--10.9 regard the SCF iteration as an approximate second-order (Newton) method, and §10.9.6 compares their convergence rates. Next to the SCF procedure on p.~64 of “qc” is written “accelerating convergence?”---this chapter answers that question.

\section{SCF Is a Fixed-Point Iteration}

One SCF iteration is a map:
\begin{align*}
D_{n+1} &= \Phi(D_n): \quad D_n \rightarrow F[D_n] \rightarrow \text{diagonalize} \rightarrow C_{n+1} \rightarrow D_{n+1}
\end{align*}
The HF solution is a fixed point $D^* = \Phi(D^*)$.

\textbf{Where does oscillation come from?} Roughly, the response of $\Phi$ is ``the Fock matrix changes a little $\to$ how much virtual character mixes into the occupied orbitals''. By first-order perturbation theory the change of an occupied orbital is proportional to $\frac{\delta F_{ai}}{\epsilon_i - \epsilon_a}$, and $\delta F$ is in turn proportional to the change in density. When the HOMO--LUMO gap is small, or electrons can ``slosh'' between several equivalent positions (large conjugated systems, transition metals), the response exceeds 1 and has the opposite sign (too many electrons go to A $\to$ in the next step the repulsion at A is too large $\to$ all the electrons go to B ...).

For example, for H$_2$O/STO-3G with both O--H bonds stretched to 1.87~Å, started from the atomic-density guess, the plain iteration (\code{rhf}) never converges: from one iteration to the next the energy alternates between $-73.5858$ and $-73.5470$~Eh. DIIS (next section), from the same start, converges in about 14 iterations to $-74.4704439437$~Eh, the PySCF value (\code{examples/diis\_vs\_plain.py}).

\section{DIIS}

\subsection{Idea}

DIIS (direct inversion in the iterative subspace, Pulay 1980, 1982) takes a linear combination of the Fock matrices of the previous steps as the Fock matrix to diagonalize next:
\begin{align*}
\tilde{F} &= \sum_{k=1}^m c_k F_k, \quad \sum_k c_k = 1
\end{align*}
How the coefficients are chosen: assume that the error is a linear function of the parameters (true near convergence); then the error of the combination is approximately $\tilde{e} = \sum_k c_k e_k$. Minimize $|\tilde{e}|^2$ under the constraint $\sum c_k = 1$. With $B_{kl} = e_k \cdot e_l$ and a Lagrange multiplier $\lambda$ for the constraint, this minimization leads to the $(m+1) \times (m+1)$ linear system
\[
\begin{pmatrix} B_{11} & \cdots & B_{1m} & -1\\ \vdots & & \vdots & \vdots\\ B_{m1} & \cdots & B_{mm} & -1\\ -1 & \cdots & -1 & 0 \end{pmatrix}
\begin{pmatrix} c_1\\ \vdots\\ c_m\\ \lambda \end{pmatrix}
= \begin{pmatrix} 0\\ \vdots\\ 0\\ -1 \end{pmatrix}
\]
Once $c$ is known, $\lambda = |\tilde{e}|^2$ is exactly the squared error after extrapolation.

\subsection{Error Vectors and Extrapolated Quantities}

In SCF the error vector is the orbital gradient $e_k = F_k D_k S - S D_k F_k$ (it vanishes exactly at the HF solution, where $f_{ai} = 0$), and the extrapolated quantity is the Fock matrix $F_k = F[D_k]$---note that both are built from the same $D_k$.

Implementation: \code{scf\_diis} in \code{hf\_to\_cc/scf.py}.

\begin{codemap}
\begin{lstlisting}[frame=none, numbers=none, xleftmargin=0pt, backgroundcolor={}]
Fs.append(F)
G = F @ d_o @ S - S @ d_o @ F
Gs.append(G)
Fs = Fs[-nkeep:]
Gs = Gs[-nkeep:]

m = len(Gs)
A = np.zeros((m+1, m+1))
for k in range(m):
    for l in range(m):
        A[k, l] = np.sum(Gs[k] * Gs[l])
A[:m, m] = -1
A[m, :m] = -1
rhs = np.zeros(m+1)
rhs[m] = -1
c = np.linalg.solve(A, rhs)
c = c[:m]

F_mix = np.einsum("k, kmn->mn", c, Fs)
\end{lstlisting}
\code{G} is the error vector $e_k = F_k D_k S - S D_k F_k$; only the last \code{nkeep} pairs are kept. \code{A} and \code{rhs} are the bordered matrix and right-hand side of the linear system above, \code{c[:m]} drops $\lambda$, and \code{F\_mix} is $\tilde{F} = \sum_k c_k F_k$, which is diagonalized in place of $F$.
\end{codemap}

\begin{pitfall}
\textbf{The initial guess decides which solution is found.} SCF only finds a stationary point of the energy, and DIIS converges quickly to a nearby one, which need not be the lowest. For N$_2$/STO-3G (\code{molecules/n2.xyz}), DIIS started from the core Hamiltonian converges in 9 iterations to $E = -106.8113762805$~Eh, while the RHF ground state is $-107.5006033119$~Eh; PySCF started from the same core-Hamiltonian guess finds the same higher solution. In that solution one virtual orbital has a negative energy ($-0.0665$~Eh) and the HOMO--LUMO gap is only $0.115$~Eh, and CCD iterated on top of it diverges. Started from a superposition of atomic densities (\code{guess\_density} in \code{provided.py}, passed to \code{scf\_diis} as \code{d\_x}), DIIS reaches $-107.5006033119$~Eh in 6 iterations.
\end{pitfall}

\begin{pitfall}
\textbf{Orbital energies come from the Fock matrix of the converged density, not from $\tilde{F}$.} Near convergence all stored $F_k$ are nearly equal and the eigenvalues of $\tilde{F}$ are the orbital energies. If the density converges within a couple of iterations, the stored matrices still include $F$ built from the initial guess. For H$_2$/STO-3G from the atomic-density guess (converged in 2 iterations, the orbitals being fixed by symmetry) the eigenvalues of $\tilde{F}$ are $(-0.5026, 0.5714)$ instead of $(-0.5797, 0.6741)$~Eh, and the MP2 energy becomes $-0.0152613$ instead of $-0.0130721$~Eh; the SCF energy itself is unaffected, because it is evaluated with the unextrapolated $F$. \code{scf\_diis} therefore diagonalizes the Fock matrix of the converged density once more before returning.
\end{pitfall}

\begin{summary}
\begin{itemize}
\item SCF is a fixed-point iteration; a small gap causes oscillations.
\item DIIS: minimize the extrapolated error under the constraint $\sum c = 1$; the error vector is the orbital gradient $\text{FDS} - \text{SDF}$.
\end{itemize}
\end{summary}

\chapter{Unrestricted Hartree--Fock}
\label{ch:uhf}

Open shells (radicals, cations, triplets) and closed-shell molecules in the bond-breaking region both require relaxing the restrictions of RHF. This chapter derives the open-shell HF method UHF, which uses a separate set of spatial orbitals for $\alpha$ and $\beta$ electrons.

\textbf{Parallel reading.} “qc” pp.~66--67 (UHF and RHF). Textbooks: Szabo--Ostlund §3.8 (the Pople--Nesbet equations of UHF; §3.8.7 is the dissociation of H$_2$ and its UHF solution).

\section{UHF: the Pople--Nesbet Equations}

\subsection{Energy}

Take spin orbitals $\chi_i^\alpha = \psi_i^\alpha \alpha$ and $\chi_i^\beta = \psi_i^\beta \beta$, with two independent sets of spatial orbitals. Start from the single-determinant energy of Chapter 2 and split the sums by spin. The one-body part splits directly:
\begin{align*}
\sum_m h_{mm} &= \sum_i^{N_\alpha} h_{ii}^\alpha + \sum_i^{N_\beta} h_{ii}^\beta
\end{align*}
Two-body part $\frac{1}{2} \sum_{mn} [\phys{mn}{mn} - \phys{mn}{nm}]$: the Coulomb term exists for all four spin combinations; the exchange term is non-zero only for equal spins (the spin-integration rules of Chapter 2):
\begin{align*}
\frac{1}{2} \sum_{mn} \phys{mn}{mn} &= \frac{1}{2} \sum_{\sigma \tau} \sum_{i \in \sigma} \sum_{j \in \tau} \chem{ii}{jj}
\end{align*}
\begin{align*}
\frac{1}{2} \sum_{mn} \phys{mn}{nm} &= \frac{1}{2} \sum_{\sigma} \sum_{i,j \in \sigma} \chem{ij}{ji}
\end{align*}
In the AO basis: $D^\sigma_{\mu \nu} = \sum_{i \in \sigma} C^\sigma_{\mu i} C^\sigma_{\nu i}$ (without the factor 2), $D = D^\alpha + D^\beta$. In the Coulomb term the sums over $i \in \sigma$ and $j \in \tau$ give $D^{\sigma}$ and $D^\tau$, and the sum over $\sigma, \tau$ gives the total density; in the exchange term both indices belong to the same spin:
\begin{align*}
\frac{1}{2} \sum_{\sigma \tau} \sum_{i \in \sigma} \sum_{j \in \tau} \chem{ii}{jj} &= \frac{1}{2} \sum_{\mu \nu \lambda \kappa} D_{\mu \nu} D_{\lambda \kappa} \chem{\mu \nu}{\lambda \kappa} = \frac{1}{2} \tr(J[D] D), \\
\frac{1}{2} \sum_{\sigma} \sum_{i,j \in \sigma} \chem{ij}{ji} &= \frac{1}{2} \sum_{\sigma} \tr(K[D^\sigma] D^\sigma)
\end{align*}
Hence
\begin{align}\label{eq:euhf}
E_{\text{UHF}} &= \tr(H D) + \frac{1}{2} \tr(J[D] D) - \frac{1}{2} \sum_{\sigma} \tr(K[D^{\sigma}] D^{\sigma}) + E_{nn}
\end{align}
(Comparison with RHF: there $D^\alpha = D^\beta = D/2$ and the exchange term is $\frac{1}{2} \cdot 2 \cdot \tr(K[D/2] D/2) = \frac{1}{4} \tr(K[D] D)$, which recovers the formula of Chapter 4.)

\subsection{Fock Matrices}

The derivative of the energy with respect to $D^{\sigma}$ is the Fock matrix. Differentiate term by term, using that $J$ and $K$ are linear in $D$ and $\tr(J[A] B) = \tr(J[B] A)$ (symmetry of the integrals). For the Coulomb term the first-order change is
\begin{align*}
\delta \left[ \frac{1}{2} \tr(J[D] D) \right] &= \frac{1}{2} \tr(J[\delta D] D) + \frac{1}{2} \tr(J[D] \delta D) \\
&= \tr(J[D] \delta D)
\end{align*}
The two terms are equal by symmetry, which cancels the $\frac{1}{2}$; and since $\delta D = \delta D^\alpha + \delta D^\beta$, the derivative with respect to $D^{\sigma}$ is $J[D]$. The exchange term is analogous but contains only the density of the same spin:
\begin{align*}
\frac{\partial}{\partial D^\sigma_{\nu \mu}} \left[ \frac{1}{2} \tr(J[D] D) \right] &= J[D]_{\mu \nu}, \\
\frac{\partial}{\partial D^\sigma_{\nu \mu}} \left[ \frac{1}{2} \tr(K[D^\sigma] D^\sigma) \right] &= K[D^\sigma]_{\mu \nu}
\end{align*}
Therefore
\begin{align}\label{eq:fock-uhf}
F^\sigma &= H + J[D^\alpha + D^\beta] - K[D^\sigma]
\end{align}
Varying each set of orbitals under the orthonormality constraint gives the Pople--Nesbet equations:
\begin{align*}
F^\alpha C^\alpha &= S C^\alpha \epsilon^\alpha, \quad F^\beta C^\beta = S C^\beta \epsilon^\beta
\end{align*}
The two equations are coupled through $J[D^\alpha + D^\beta]$. The energy can be written $E = \frac{1}{2} \sum_{\sigma} \tr[(H + F^{\sigma}) D^{\sigma}] + E_{nn}$.

Implementation: \code{uhf} (plain iteration) and \code{uhf\_diis} (with the DIIS of Chapter 5) in \code{hf\_to\_cc/scf.py}.

\begin{codemap}
\begin{lstlisting}[frame=none, numbers=none, xleftmargin=0pt, backgroundcolor={}]
j = np.einsum("mnls, ls->mn", eri, da_o + db_o)
ka = np.einsum("mlns, ls->mn", eri, da_o)
kb = np.einsum("mlns, ls->mn", eri, db_o)
fa = h + j - ka
fb = h + j - kb
\end{lstlisting}
\eqref{eq:fock-uhf} for $\sigma = \alpha, \beta$: one Coulomb matrix $J[D^\alpha + D^\beta]$ from the total density, one exchange matrix $K[D^\sigma]$ per spin from the density of that spin (\code{da\_o}, \code{db\_o} are $D^\alpha$, $D^\beta$, without the factor 2).
\end{codemap}

\begin{pitfall}
\textbf{UHF has several solutions.} Converging is not enough. For CO$_2^+$/STO-3G (\code{molecules/co2.xyz}, charge $+1$), \code{uhf\_diis} started from the core Hamiltonian converges to $-184.3961$~Eh, and started from the orbitals of the atomic-density guess (\code{guess\_orbitals}) to $-184.7402$~Eh, the PySCF value. For HCN$^+$/STO-3G it is the other way round: the core-Hamiltonian start gives a solution 17~mEh \emph{below} PySCF's. For N$_2^+$ every start ends on a solution that is not a minimum. Which solution is the lowest can only be decided by examining the second derivatives of the energy with respect to orbital rotations (a stability analysis).
\end{pitfall}

\begin{summary}
\begin{itemize}
\item UHF: $F^{\sigma} = H + J[D^{\alpha} + D^{\beta}] - K[D^{\sigma}]$.
\end{itemize}
\end{summary}

\chapter{Second Quantization and Wick's Theorem}
\label{ch:sq}

The coupled-cluster derivations of Chapter 11 are all done in second quantization: it ``packs'' the antisymmetry of determinants into the anticommutation relations of operators, and the Slater--Condon rules become simple algebra. This chapter follows the order of MEST pp.~21--27: creation and annihilation operators, anticommutation relations, operators in second quantization, normal order and Wick's theorem, the Fermi vacuum and the particle--hole formalism, and the normal-ordered Hamiltonian. Finally it explains how $\cre{p}$ acts on a determinant stored as a bit string.

\textbf{Parallel reading.} “Adv qc” pp.~15--20 (creation and annihilation operators, anticommutation relations, one- and two-body operators) and pp.~40--47 (normal order, contractions, Wick's theorem, the Fermi vacuum, the normal-ordered Hamiltonian, the generalized Wick theorem); MEST pp.~21--26, 40--41. The textbook closest in order to this chapter is Shavitt--Bartlett Ch.~3 (§3.2 creation and annihilation operators, §3.3 normal products and Wick's theorem, §3.4 the particle--hole formalism, §3.6 operators in normal-product form, §3.7 the generalized Wick theorem, §3.8 evaluation of matrix elements); Helgaker Ch.~1--2 and ESQC collection Ch.~6--12 (Olsen) do not use Wick's theorem but commutator algebra and singlet excitation operators, which makes a good comparison with this chapter; Szabo--Ostlund §2.4 is a short introduction. Exercises: ESQC 76.1 (six identities for commutators and anticommutators), 76.8 (commutation relations between $E_{mn}$ and creation/annihilation operators); the introduction of ESQC 2024 problem 80 restates the Fermi-vacuum formalism in terms of quasiparticle operators.

\section{Creation and Annihilation Operators}

\subsection{Definitions}

Take an orthonormal set of spin orbitals $\no{\chi_p}$. Define the creation operator $\cre{p}$:
\begin{align*}
\cre{p} |\chi_k \chi_l \dots \chi_n\rangle = |\chi_p \chi_k \chi_l \dots \chi_n\rangle
\end{align*}
i.e.\ it puts $\chi_p$ at the \emph{front} of the determinant. If $\chi_p$ is already occupied, the result is a determinant with two equal columns, which is zero. Every determinant can be generated from the vacuum $|vac\rangle$ (no electrons, normalized $\langle vac|vac\rangle = 1$):
\begin{align*}
|\chi_i \chi_j \dots \chi_k\rangle = \cre{i} \cre{j} \dots \cre{k} |\text{vac}\rangle
\end{align*}
The annihilation operator is $\ann{p} = (\cre{p})^\dagger$; acting on a determinant with $\chi_p$ at the front it gives
\begin{align*}
a_p |\chi_p \chi_i \dots\rangle &= |\chi_i \dots\rangle
\end{align*}
i.e.\ $\ann{p}$ removes $\chi_p$ \emph{standing at the front}. If $\chi_p$ is not at the front, first swap columns to bring it there (a minus sign per swap), then remove it. If $\chi_p$ is not occupied, $\ann{p}$ gives zero. The vacuum has nothing to annihilate: $\ann{p} |vac\rangle = 0$, and correspondingly $\langle vac| \cre{p} = 0$.

\textbf{Fock space.} $\cre{p}$ raises the number of electrons by one and $\ann{p}$ lowers it by one, so the space they act on contains states of every electron number: the vacuum, the one-electron states $\cre{p} |\text{vac}\rangle$, the two-electron states $\cre{p} \cre{q} |\text{vac}\rangle$, and so on. These fixed-number spaces of antisymmetric states taken together form the Fock space. Every term of the electronic Hamiltonian below contains as many creation as annihilation operators and conserves the number of electrons, so CI is solved within the $N$-electron part alone. With $n$ spin orbitals, each occupied or empty, the Fock space has dimension $2^n$.

\subsection{Anticommutation Relations}

Comparing $\cre{p} \cre{q}$ with $\cre{q} \cre{p}$, and $\ann{p} \cre{q}$ with $\cre{q} \ann{p}$, on every determinant gives the anticommutation relations ($\{A, B\} = A B + B A$)
\begin{align}\label{eq:anticomm}
\{\cre{p}, \cre{q}\} = 0, \quad \{a_p, a_q\} = 0, \quad \{a_p, \cre{q}\} = \delta_{pq}
\end{align}

\section{Operators in Second Quantization}

The one-body operator $O_1 = \sum_i h(i)$ and the two-body operator $O_2 = \sum_{i<j} v(i,j)$ are written in second quantization as
\begin{align*}
O_1 = \sum_{pq} \langle p|h|q\rangle \cre{p} a_q
\end{align*}
\begin{align}\label{eq:op2}
&O_2 = \frac{1}{2} \sum_{pqrs} \phys{pq}{rs} \cre{p} \cre{q} a_s a_r \\
&= \frac{1}{4} \sum_{pqrs} \asym{pq}{rs} \cre{p} \cre{q} a_s a_r \notag
\end{align}
Note that the order of the annihilation operators in $O_2$ is $\ann{s} \ann{r}$ (opposite to the order of $r, s$ in the integral $\phys{pq}{rs}$).

\textbf{Check $\frac{1}{2} \phys{pq}{rs} \rightarrow \frac{1}{4} \asym{pq}{rs}$:} in $\frac{1}{2} \sum \phys{pq}{rs} \cre{p} \cre{q} \ann{s} \ann{r}$ swap the dummy indices $r \leftrightarrow s$ and use $\ann{r} \ann{s} = - \ann{s} \ann{r}$:
\begin{align*}
\frac{1}{2} \sum \phys{pq}{rs} \cre{p} \cre{q} a_s a_r = - \frac{1}{2} \sum \phys{pq}{sr} \cre{p} \cre{q} a_s a_r
\end{align*}
Adding the two forms and dividing by 2 gives $\frac{1}{4} \sum \asym{pq}{rs} \cre{p} \cre{q} \ann{s} \ann{r}$.

\section{Normal Order, Contractions and Wick's Theorem}

\subsection{Normal Order with Respect to the Vacuum}

The normal-ordered product $\{ABC\dots\}$ of a string of operators: move all annihilation operators to the right of all creation operators, with a factor $-1$ for every swap (\emph{without} using the $\delta$ terms of the anticommutation relations---it is only a reordering). For example
\begin{align*}
\no{\cre{p} a_q} &= \cre{p} a_q, \quad \no{a_q \cre{p}} = - \cre{p} a_q, \quad \no{a_p a_q} = a_p a_q
\end{align*}
The vacuum expectation value of a normal-ordered product is zero (the annihilation operators on the right act on $|vac\rangle$, or the creation operators on the left act on $\langle \text{vac}|$).

\subsection{Contractions}

The contraction of two operators is defined as the difference between their product and its normal-ordered product:
\begin{align*}
\contraction{A B} = A B - \no{A B}
\end{align*}
With respect to the vacuum, the only non-zero contraction is
\begin{align*}
\contraction{a_p \cre{q}} = a_p \cre{q} - \no{a_p \cre{q}} = a_p \cre{q} + \cre{q} a_p = \delta_{pq}
\end{align*}
The other three kinds ($\ann{}^\dagger \ann{}^\dagger$, $\ann{} \ann{}$, $\ann{}^\dagger \ann{}$) vanish because they are already normal ordered.

\subsection{Wick's Theorem}

\begin{theorem}[Wick]
A product of creation and annihilation operators equals its normal-ordered product, plus the normal-ordered products with all single contractions, plus all double contractions, \ldots, up to all fully contracted terms:
\begin{align*}
A B C \dots X Y Z &= \no{A B C \dots X Y Z} + \sum_{\text{single}} \no{\contraction \dots} \\
&\quad + \sum_{\text{double}} \no{\dots} + \dots + \sum_{\text{full}} \no{\text{fully contracted}}
\end{align*}
The sign of each term is the parity of a reordering: bring each contracted pair of operators next to each other in turn (the uncontracted operators keep their relative order); an odd number of swaps gives a minus sign.
\end{theorem}

\textbf{Counting crossings for full contractions.} When all operators are contracted, this parity is $(-1)^h$, where $h$ is the number of crossings of the contraction lines; this is the form used for matrix elements from here on.

The value of Wick's theorem: every matrix element equals the sum of the \emph{fully contracted} terms, and full contractions only require listing the pairings and counting the crossings.

\section{The Fermi Vacuum and the Particle--Hole Formalism}

The vacuum $|vac\rangle$ is too far from the physics: we care about $N$-electron states, all of which are ``close to'' the HF determinant $|\Phi_0\rangle$. Take $|\Phi_0\rangle$ as the new ``vacuum''---the Fermi vacuum. Following the convention of MEST p.~25, the occupied orbitals $i, j, k, l$ are called hole orbitals and the virtual orbitals $a, b, c, d$ particle orbitals. Define quasiparticle operators:
\begin{align*}
&\text{hole creation operator: } a_i \ \text{(removes an electron from } \Phi_0 = \text{creates a hole)}, \\
&\text{hole annihilation operator: } \cre{i}
\end{align*}
\begin{align*}
&\text{particle creation operator: } \cre{a}, \quad \text{particle annihilation operator: } a_a
\end{align*}
Because $\cre{i} |\Phi_0\rangle = 0$ ($i$ is occupied) and $\ann{a} |\Phi_0\rangle = 0$ ($a$ is empty), the quasiparticle annihilation operators give zero on the Fermi vacuum, just as in the vacuum case.

\textbf{Normal order $\{ \dots \}$ with respect to $|\Phi_0\rangle$}: move all quasiparticle annihilation operators ($\cre{i}$, $\ann{a}$) to the right. Only two \textbf{contractions} are non-zero:
\begin{align*}
\contraction{\cre{i} a_j} = \delta_{ij}, \quad \contraction{a_a \cre{b}} = \delta_{ab}
\end{align*}
Derivation: for $\cre{i} \ann{j} - \{\cre{i} \ann{j}\}$, $\cre{i}$ is a hole annihilation operator and must move to the right, $\{\cre{i} \ann{j}\} = -\ann{j} \cre{i}$, so the contraction is $\cre{i} \ann{j} + \ann{j} \cre{i} = \delta_{ij}$. $\ann{a} \cre{b}$ is the same as in the vacuum case. All other combinations ($\ann{i} \cre{j}$, $\cre{a} \ann{b}$, mixed $i a$) have zero contraction. Mnemonic: \emph{hole lines contract in the order $\ann{}^\dagger \ann{}$, particle lines in the order $a\, a^\dagger$.}

With these rules the \emph{expectation value} of a normal-ordered product with respect to $|\Phi_0\rangle$ is zero, and Wick's theorem holds unchanged. Note that a normal-ordered product acting on $|\Phi_0\rangle$ need not give zero: for example $\no{\cre{a} a_i}|\Phi_0\rangle = |\Phi_i^a\rangle$; only products whose rightmost operator is a quasiparticle annihilator ($\cre{i}$ or $\ann{a}$) annihilate $|\Phi_0\rangle$.

\section{The Normal-Ordered Hamiltonian}

Both parts of $H$ are expanded with Wick's theorem, with normal order and contractions taken with respect to the Fermi vacuum $|\Phi_0\rangle$ (MEST p.~26; “Adv qc” pp.~43--45):
\begin{align*}
H &= O_1 + O_2 = \sum_{pq} h_{pq} \cre{p} \ann{q} + \frac{1}{2} \sum_{pqrs} \phys{pq}{rs} \cre{p} \cre{q} \ann{s} \ann{r}
\end{align*}

\textbf{One-electron part.} Wick's theorem gives one term without and one with a contraction:
\begin{align*}
O_1 &= \sum_{pq} h_{pq} \no{\cre{p} \ann{q}} + \sum_{pq} h_{pq} \no{\contraction{\cre{p} \ann{q}}} \\
&= \sum_{pq} h_{pq} \no{\cre{p} \ann{q}} + \sum_{ij} h_{ij} \contraction{\cre{i} \ann{j}} \eqnote{hole line: only occupied $p, q$} \\
&= \sum_{pq} h_{pq} \no{\cre{p} \ann{q}} + \sum_{i} h_{ii} \eqnote{$\contraction{\cre{i} \ann{j}} = \delta_{ij}$}
\end{align*}

\textbf{Two-electron part.} In $\cre{p} \cre{q} \ann{s} \ann{r}$ two $a^\dagger$ or two $a$ never contract, and $\cre{p}$ standing to the left of $\ann{s}$ gives $\contraction{\cre{p} \ann{s}} = \delta^{\text{occ}}_{ps}$, which is $1$ if $p = s$ is occupied and $0$ otherwise. This leaves four single and two double contractions (two crossing lines are drawn one above and one below the string):
% crossing pair p--s, q--r: p--s above (\contraction over a phantom), q--r below (\underbracket)
\begin{align*}
\cre{p} \cre{q} \ann{s} \ann{r} &= \no{\cre{p} \cre{q} \ann{s} \ann{r}} + \no{{\contraction{\cre{p} \cre{q} \ann{s}}} \ann{r}} + \no{\contraction{\cre{p} \cre{q} \ann{s} \ann{r}}} \\
&\quad + \no{\cre{p} {\contraction{\cre{q} \ann{s}}} \ann{r}} + \no{\cre{p} {\contraction{\cre{q} \ann{s} \ann{r}}}} \\
&\quad + \no{\mathrlap{\contraction{\phantom{\cre{p} \cre{q} \ann{s}}}} \cre{p} {\underbracket{\cre{q} \ann{s} \ann{r}}}} + \no{\contraction{\cre{p} {\contraction{\cre{q} \ann{s}}} \ann{r}}}
\end{align*}
The sign of each term is the parity of the swaps that bring the contracted operators together:
\begin{align*}
\cre{p} \cre{q} \ann{s} \ann{r} &= \no{\cre{p} \cre{q} \ann{s} \ann{r}} - \delta^{\text{occ}}_{ps} \no{\cre{q} \ann{r}} + \delta^{\text{occ}}_{pr} \no{\cre{q} \ann{s}} \eqnote{1 and 2 swaps} \\
&\quad + \delta^{\text{occ}}_{qs} \no{\cre{p} \ann{r}} - \delta^{\text{occ}}_{qr} \no{\cre{p} \ann{s}} \eqnote{0 and 1 swaps} \\
&\quad - \delta^{\text{occ}}_{ps} \delta^{\text{occ}}_{qr} + \delta^{\text{occ}}_{pr} \delta^{\text{occ}}_{qs} \eqnote{crossing and nested}
\end{align*}
Multiply by $\frac{1}{2} \phys{pq}{rs}$ and sum over $p, q, r, s$; each $\delta^{\text{occ}}$ leaves one occupied index ($i$, and $j$ for the second one):
\begin{align*}
O_2 &= \frac{1}{2} \sum_{pqrs} \phys{pq}{rs} \no{\cre{p} \cre{q} \ann{s} \ann{r}} - \frac{1}{2} {\underbrace{\sum_{iqr} \phys{iq}{ri} \no{\cre{q} \ann{r}}}_{(1)}} \\
&\quad + \frac{1}{2} {\underbrace{\sum_{iqs} \phys{iq}{is} \no{\cre{q} \ann{s}}}_{(3)}} + \frac{1}{2} {\underbrace{\sum_{pir} \phys{pi}{ri} \no{\cre{p} \ann{r}}}_{(4)}} \\
&\quad - \frac{1}{2} {\underbrace{\sum_{pis} \phys{pi}{is} \no{\cre{p} \ann{s}}}_{(2)}} - \frac{1}{2} \sum_{ij} \phys{ij}{ji} + \frac{1}{2} \sum_{ij} \phys{ij}{ij}
\end{align*}
The four single-contraction sums are equal in pairs (first rename the summation indices, then exchange the two integration variables):
\begin{align*}
(2) &= \sum_{iqr} \phys{qi}{ir} \no{\cre{q} \ann{r}} = \sum_{iqr} \phys{iq}{ri} \no{\cre{q} \ann{r}} = (1) \eqnote{$p \to q$, $s \to r$} \\
(4) &= \sum_{iqs} \phys{qi}{si} \no{\cre{q} \ann{s}} = \sum_{iqs} \phys{iq}{is} \no{\cre{q} \ann{s}} = (3) \eqnote{$p \to q$, $r \to s$}
\end{align*}
Hence, with $\asym{ij}{ij} = \phys{ij}{ij} - \phys{ij}{ji}$ and, in the second step, the summation indices of the two single-contraction sums renamed to $p, q$,
\begin{align*}
O_2 &= \frac{1}{2} \Big[ \sum_{pqrs} \phys{pq}{rs} \no{\cre{p} \cre{q} \ann{s} \ann{r}} - 2 \sum_{iqr} \phys{qi}{ir} \no{\cre{q} \ann{r}} + 2 \sum_{iqs} \phys{qi}{si} \no{\cre{q} \ann{s}} + \sum_{ij} \asym{ij}{ij} \Big] \\
&= \frac{1}{2} \Big[ \sum_{pqrs} \phys{pq}{rs} \no{\cre{p} \cre{q} \ann{s} \ann{r}} + 2 \sum_{pqi} \phys{pi}{qi} \no{\cre{p} \ann{q}} - 2 \sum_{pqi} \phys{pi}{iq} \no{\cre{p} \ann{q}} + \sum_{ij} \asym{ij}{ij} \Big] \\
&= \frac{1}{2} \sum_{pqrs} \phys{pq}{rs} \no{\cre{p} \cre{q} \ann{s} \ann{r}} + \sum_{pqi} \asym{pi}{qi} \no{\cre{p} \ann{q}} + \frac{1}{2} \sum_{ij} \asym{ij}{ij}
\end{align*}

\textbf{Collecting terms.} Adding $O_1$:
\begin{align*}
H &= \sum_{pq} \Big[ h_{pq} + \sum_i \asym{pi}{qi} \Big] \no{\cre{p} \ann{q}} + \frac{1}{2} \sum_{pqrs} \phys{pq}{rs} \no{\cre{p} \cre{q} \ann{s} \ann{r}} \\
&\quad + \sum_i h_{ii} + \frac{1}{2} \sum_{ij} \asym{ij}{ij}
\end{align*}
The bracket is the Fock matrix, $f_{pq} = h_{pq} + \sum_i \asym{pi}{qi} = \langle p|\hat{f}|q\rangle$ with $\hat{f}$ of \eqref{eq:fock-op}. In the two-electron term, the swap $r \leftrightarrow s$ of \eqref{eq:op2} turns $\frac{1}{2} \phys{pq}{rs}$ into $\frac{1}{4} \asym{pq}{rs}$ (inside $\no{\dots}$ a swap of two operators also gives $-1$). The constant is $\langle \Phi_0|H|\Phi_0\rangle$, since every normal-ordered term has zero expectation value, and by \eqref{eq:E-det} it is $E_{\text{HF}}$:
\begin{align}
H &= \sum_{pq} f_{pq} \no{\cre{p} \ann{q}} + \frac{1}{4} \sum_{pqrs} \asym{pq}{rs} \no{\cre{p} \cre{q} \ann{s} \ann{r}} + \langle \Phi_0|H|\Phi_0\rangle \notag \\
&= \FN + \VN + E_{\text{HF}} \label{eq:HN}
\end{align}
$\FN$ is the Fock operator and $\VN$ the two-electron operator, both in normal-ordered form. The normal-ordered Hamiltonian
\begin{align*}
\HN &= H - E_{\text{HF}} = \FN + \VN
\end{align*}
is the starting point of MBPT and CC.

\begin{remark}
$f_{pq} = h_{pq} + \sum_i \asym{pi}{qi}$ holds for \emph{any} reference determinant, as long as $i$ runs over the occupied orbitals of that determinant.
\end{remark}

\section{The Generalized Wick Theorem and Matrix Elements}

\textbf{Generalized Wick theorem} (MEST p.~26; “Adv qc” p.~46). For a product of strings that are each already normal ordered,
\begin{align*}
\no{A B C \cdots} \no{X Y Z \cdots} &= \no{A B C \cdots X Y Z \cdots} + \sum_{\text{single}} \no{{\contraction{A B C \cdots X}} Y Z \cdots} \\
&\quad + \sum_{\text{double}} \no{\mathrlap{\contraction{\phantom{A B C \cdots X}}} A {\underbracket{B C \cdots X Y}} Z \cdots} + \cdots
\end{align*}
where contractions are allowed only between operators of \emph{different} normal-ordered strings. As in Wick's theorem, $\langle \Phi_0|\no{A B C \cdots} \no{X Y Z \cdots}|\Phi_0\rangle$ is non-zero only for the fully contracted terms, in which every operator is paired into a contraction; so the two strings must contain equally many operators. The excited determinants are such strings acting on $|\Phi_0\rangle$ ($\cre{a}$ and $\ann{i}$ are quasiparticle creation operators):
\begin{align*}
|\Phi_i^a\rangle &= \cre{a} \ann{i} |\Phi_0\rangle = \no{\cre{a} \ann{i}} |\Phi_0\rangle \\
|\Phi_{ij}^{ab}\rangle &= \cre{a} \cre{b} \ann{j} \ann{i} |\Phi_0\rangle = \no{\cre{a} \cre{b} \ann{j} \ann{i}} |\Phi_0\rangle
\end{align*}

\textbf{Single excitation} (MEST p.~27; “Adv qc” pp.~46--47). With \eqref{eq:HN},
\begin{align*}
\langle \Phi_0|H|\Phi_i^a\rangle &= \langle \Phi_0|H \cre{a} \ann{i}|\Phi_0\rangle \\
&= \langle \Phi_0|(\FN + \VN + E_{\text{HF}}) \no{\cre{a} \ann{i}}|\Phi_0\rangle
\end{align*}
where
\begin{align*}
\langle \Phi_0|E_{\text{HF}} \no{\cre{a} \ann{i}}|\Phi_0\rangle &= E_{\text{HF}} \langle \Phi_0|\Phi_i^a\rangle = 0 \\
\langle \Phi_0|\VN \no{\cre{a} \ann{i}}|\Phi_0\rangle &= 0 \eqnote{4 operators cannot all pair with 2}
\end{align*}
In the $\FN$ term $\cre{p}$ contracts with $\ann{i}$ (hole line) and $\ann{q}$ with $\cre{a}$ (particle line); the two lines are nested, so the sign is $+$:
\begin{align*}
\langle \Phi_0|H|\Phi_i^a\rangle &= \sum_{pq} f_{pq} \langle \Phi_0|\no{\cre{p} \ann{q}} \no{\cre{a} \ann{i}}|\Phi_0\rangle \\
&= \sum_{pq} f_{pq} \langle \Phi_0|\no{\contraction{\cre{p} {\contraction{\ann{q} \cre{a}}} \ann{i}}}|\Phi_0\rangle \\
&= \sum_{pq} f_{pq} \delta_{pi} \delta_{qa} \langle \Phi_0|\Phi_0\rangle \\
&= f_{ia}
\end{align*}
For HF orbitals $f_{ia} = 0$: Brillouin's theorem (Chapter~3).

\textbf{Double excitation} (MEST p.~27; “Adv qc” p.~47). Now the $E_{\text{HF}}$ and $\FN$ terms vanish (two operators cannot all pair with four), and $\VN$ is used in the form $\frac{1}{2} \sum \phys{pq}{rs} \no{\cre{p} \cre{q} \ann{s} \ann{r}}$:
\begin{align*}
\langle \Phi_0|H|\Phi_{ij}^{ab}\rangle &= \langle \Phi_0|H \cre{a} \cre{b} \ann{j} \ann{i}|\Phi_0\rangle \\
&= \frac{1}{2} \sum_{pqrs} \phys{pq}{rs} \langle \Phi_0|\no{\cre{p} \cre{q} \ann{s} \ann{r}} \no{\cre{a} \cre{b} \ann{j} \ann{i}}|\Phi_0\rangle
\end{align*}
Number the operators $\cre{p} (1)$, $\cre{q} (2)$, $\ann{s} (3)$, $\ann{r} (4)$, $\cre{a} (5)$, $\cre{b} (6)$, $\ann{j} (7)$, $\ann{i} (8)$. In a full contraction $\ann{s}, \ann{r}$ pair with $\cre{a}, \cre{b}$ (particle lines) and $\cre{p}, \cre{q}$ with $\ann{j}, \ann{i}$ (hole lines). This gives $2 \times 2$ terms, each with the sign $(-1)^h$ for $h$ crossings:
\begin{align*}
\langle \Phi_0|\no{\cre{p} \cre{q} \ann{s} \ann{r}} \no{\cre{a} \cre{b} \ann{j} \ann{i}}|\Phi_0\rangle &= \delta_{sa} \delta_{rb} \delta_{pj} \delta_{qi} \eqnote{lines 3--5, 4--6, 1--7, 2--8: $h = 2$} \\
&\quad - \delta_{sb} \delta_{ra} \delta_{pj} \delta_{qi} \eqnote{lines 3--6, 4--5, 1--7, 2--8: $h = 1$} \\
&\quad - \delta_{sa} \delta_{rb} \delta_{pi} \delta_{qj} \eqnote{lines 3--5, 4--6, 1--8, 2--7: $h = 1$} \\
&\quad + \delta_{sb} \delta_{ra} \delta_{pi} \delta_{qj} \eqnote{lines 3--6, 4--5, 1--8, 2--7: $h = 0$}
\end{align*}
Hence
\begin{align*}
\langle \Phi_0|H|\Phi_{ij}^{ab}\rangle &= \frac{1}{2} \big[ \phys{ji}{ba} - \phys{ji}{ab} - \phys{ij}{ba} + \phys{ij}{ab} \big] \\
&= \frac{1}{2} \big[ 2 \phys{ij}{ab} - 2 \phys{ij}{ba} \big] \eqnote{exchange the two integration variables} \\
&= \asym{ij}{ab}
\end{align*}
Both results agree with the Slater--Condon rules of Chapter~2 (cases 2 and 3). The bra forms used in Chapters 9 and 11 are the complex conjugates ($H$ is Hermitian), with $\langle \Phi_{ij}^{ab}| = \langle \Phi_0|\cre{i} \cre{j} \ann{b} \ann{a} = \langle \Phi_0|\no{\cre{i} \cre{j} \ann{b} \ann{a}}$:
\begin{align*}
\langle \Phi_i^a|H|\Phi_0\rangle &= f_{ai} \\
\langle \Phi_{ij}^{ab}|\VN|\Phi_0\rangle &= \langle \Phi_{ij}^{ab}|H|\Phi_0\rangle = \asym{ab}{ij}
\end{align*}

\section{Determinants as Bit Strings}

A program that needs determinants (FCI, for example) can store each one as an integer whose bit $p$ is 1 when spin orbital (or spatial orbital) $p$ is occupied, ordered by increasing index,
\begin{align*}
|D\rangle &= \cre{p_1} \cre{p_2} \dots \cre{p_N} |\text{vac}\rangle, \quad p_1 < p_2 < \dots < p_N
\end{align*}
\textbf{Sign rule.} To get to its correct position, $\cre{q}$ has to pass the creation operators of all occupied orbitals with index below $q$ (by \eqref{eq:anticomm}, a factor $-1$ for each):
\begin{align*}
\cre{q} |D\rangle &= (-1)^{N_{<q}} |D + q\rangle, \quad a_q |D\rangle = (-1)^{N_{<q}} |D - q\rangle
\end{align*}
where $N_{<q}$ is the number of occupied orbitals with index below $q$.

\begin{summary}
\begin{itemize}
\item $\{a_p, \cre{q}\} = \delta_{pq}$ and all other anticommutators vanish; the Pauli principle and the antisymmetry of determinants are contained in them.
\item $O_2 = \frac{1}{4} \sum \asym{pq}{rs} \cre{p} \cre{q} a_s a_r$ (note the order $\ann{s} \ann{r}$).
\item With respect to the Fermi vacuum the only non-zero contractions are $\cre{i} a_j \rightarrow \delta_{ij}$ and $a_a \cre{b} \rightarrow \delta_{ab}$.
\item $H = E_{\text{HF}} + \FN + \VN$, with $f_{pq} = h_{pq} + \sum_i \asym{pi}{qi}$ valid for any reference determinant.
\item Determinant = bit string, sign $(-1)^{N_{<p}}$.
\end{itemize}
\end{summary}

\chapter{Integral Transformation and Spin-Orbital Integrals}
\label{ch:integrals}

All post-HF methods need integrals in the molecular-orbital (MO) representation. This chapter discusses two things: (1) how the AO$\to$MO transformation is reduced from $\mathcal{O}(N^8)$ to $\mathcal{O}(N^5)$; (2) the construction of the antisymmetrized spin-orbital integrals $\asym{pq}{rs}$ and the Fock matrix, the input of coupled cluster (Chapter 11).

\textbf{Parallel reading.} The Sherrill--Schaefer CI review §3.1 (transformation of one- and two-electron integrals); Python for Quantum Chemistry §12.3 (integral transformation).

\section{The AO$\to$MO Transformation}

\subsection{Four Quarter Transformations}

\begin{align*}
\chem{pq}{rs} &= \sum_{\mu \nu \lambda \sigma} C_{\mu p} C_{\nu q} \chem{\mu \nu}{\lambda \sigma} C_{\lambda r} C_{\sigma s}
\end{align*}
Writing this formula directly as one \code{einsum}, each output element sums over four AO indices: $N^4$ outputs × $N^4$ summations = $\mathcal{O}(N^8)$. The key is to transform one index at a time:
\begin{align*}
\chem{\mu \nu}{\lambda s} &= \sum_{\sigma} \chem{\mu \nu}{\lambda \sigma} C_{\sigma s} \quad \mathcal{O}(N^5)
\end{align*}
\begin{align*}
\chem{\mu \nu}{r s} &= \sum_{\lambda} C_{\lambda r} \chem{\mu \nu}{\lambda s} \quad \mathcal{O}(N^5)
\end{align*}
\begin{align*}
\chem{\mu q}{r s} &= \sum_{\nu} C_{\nu q} \chem{\mu \nu}{r s} \quad \mathcal{O}(N^5)
\end{align*}
\begin{align*}
\chem{p q}{r s} &= \sum_{\mu} C_{\mu p} \chem{\mu q}{r s} \quad \mathcal{O}(N^5)
\end{align*}
Each step sums over a single index at cost $\mathcal{O}(N^5)$, four steps in total. This is the basic idea of ``splitting a multiple sum into a sequence of pairwise contractions''. \code{einsum(..., optimize=True)} finds such contraction orders automatically, but for the integral transformation we write the four steps by hand.

Implementation: \code{ao2mo} in \code{hf\_to\_cc/ao2mo.py}.

\begin{codemap}
\begin{lstlisting}[frame=none, numbers=none, xleftmargin=0pt, backgroundcolor={}]
mo = np.einsum("ma, mnls->anls", C, eri)
mo = np.einsum("nb, mnls->mbls", C, mo)
mo = np.einsum("lc, mnls->mncs", C, mo)
mo = np.einsum("sd, mnls->mnld", C, mo)
\end{lstlisting}
Each line is one quarter transformation, a single sum over one AO index at cost $\mathcal{O}(N^5)$. The code transforms $\mu$ first and $\sigma$ last, the reverse of the order above; any order gives the same $\chem{pq}{rs}$, returned as \code{mo[p, q, r, s]}.
\end{codemap}

\section{Spin-Orbital Integrals}

\subsection{Ordering}

The spin orbitals are ordered by spatial orbital, with the two spins interleaved as in Chapter 2 (counting from 0, spin orbital $2p$ is $\psi_p \alpha$ and $2p + 1$ is $\psi_p \beta$):
\begin{align*}
&[\psi_0 \alpha, \psi_0 \beta, \psi_1 \alpha, \psi_1 \beta, \dots]
\end{align*}
For a closed shell the first $N$ are occupied. The $o$/$v$ blocking is then a simple slice \code{[:N]}, \code{[N:]}. Each spin orbital $P$ has the spatial orbital $\psi_{\lfloor P/2 \rfloor}$ and a spin label $s_P$ ($\alpha$ for even $P$, $\beta$ for odd $P$). For RHF, $\alpha$ and $\beta$ use the same set of spatial orbitals; for UHF each uses its own.

\subsection{Antisymmetrized Integrals}

First do one AO$\to$MO transformation with the spatial coefficients of all spin orbitals to obtain $\chem{PQ}{RS}_{\text{space}}$, then multiply by the spin-integration factor (Chapter 2: the two orbitals of the same electron must have the same spin):
\begin{align*}
(PQ|RS) &= (PQ|RS)_{\text{space}} \delta_{s_P s_Q} \delta_{s_R s_S}
\end{align*}
Convert to physicists' notation and antisymmetrize:
\begin{align*}
\phys{PQ}{RS} &= (PR|QS), \quad \asym{PQ}{RS} = \phys{PQ}{RS} - \phys{PQ}{SR} = (PR|QS) - (PS|QR)
\end{align*}
\code{spin\_orbitals} in \code{hf\_to\_cc/ao2mo.py} returns $\asym{PQ}{RS}$, in this ordering, as the four-index array \code{g[P, Q, R, S]}.

\begin{codemap}
\begin{lstlisting}[frame=none, numbers=none, xleftmargin=0pt, backgroundcolor={}]
v = 0.0
if p % 2 == r % 2 and q % 2 == s % 2:
    v += mo[p // 2, r // 2, q // 2, s // 2]
if p % 2 == s % 2 and q % 2 == r % 2:
    v -= mo[p // 2, s // 2, q // 2, r // 2]
g[p, q, r, s] = v
\end{lstlisting}
Spin orbital \code{p} is spatial orbital \code{p // 2} with spin \code{p \% 2} ($0 = \alpha$, $1 = \beta$). The first \code{if} is the spin condition of $\phys{PQ}{RS} = (PR|QS)$, the second that of the exchange term $(PS|QR)$.
\end{codemap}

\subsection{Fock Matrix and Reference Energy}

The Fock matrix of \eqref{eq:HN} of Chapter 7 is $f_{pq} = h_{pq} + \sum_{i \in \text{occ}} \asym{pi}{qi}$; for RHF/UHF references it is diagonal. The reference energy is
\begin{align*}
E_{\text{ref}} &= E_{nn} + \sum_i h_{ii} + \frac{1}{2} \sum_{ij} \asym{ij}{ij}
\end{align*}
$E_{\text{ref}}$ must agree with the HF energy, which makes it an overall check that the integral transformation and the spin integration are both correct.

\begin{summary}
\begin{itemize}
\item Four $\mathcal{O}(N^5)$ quarter transformations instead of one $\mathcal{O}(N^8)$ sum.
\item Spin-orbital integrals = spatial integrals × spin deltas, then $\asym{PQ}{RS} = \chem{PR}{QS} - \chem{PS}{QR}$; the Fock matrix $f = h + \sum_i \asym{pi}{qi}$ holds for any reference.
\end{itemize}
\end{summary}

\chapter{Perturbation Theory and MP2}
\label{ch:mp2}

MP2 is the cheapest correlation method and also the starting guess of the CC iterations (Chapter 11). This chapter follows the order of MEST pp.~28--32: first Rayleigh--Schrödinger perturbation theory (RSPT) is recalled, then the Møller--Plesset partitioning gives MP2; the spin-orbital formula is then spin-integrated into the closed-shell formula (every spin case written out).

\textbf{Parallel reading.} “Adv qc” pp.~48--55: from MBPT for a one-electron perturbation to the MP partitioning and MP2. Also “qc” pp.~84--86 (MP partitioning, $E_0 \neq E_{\text{HF}}$) and MEST pp.~28--32. Textbooks: Shavitt--Bartlett Ch.~2 (derivation of RS perturbation theory); Helgaker §14.1, §14.2 (MP theory), §14.4 (closed-shell MP; Eq.~(14.4.40) defines $L_{pqrs} = 2 g_{pqrs} - g_{psrq}$ and Eq.~(14.4.55) is the closed-shell MP2 energy); Szabo--Ostlund Ch.~6.

\section{Rayleigh--Schrödinger Perturbation Theory}

This chapter uses only a few results of non-degenerate RS perturbation theory. Let $\hat{H} = \hat{H}_0 + \lambda \hat{V}$ with $\hat{H}_0 |\Phi_n^{(0)}\rangle = E_n^{(0)} |\Phi_n^{(0)}\rangle$, let the target zeroth-order state $|\Phi_i^{(0)}\rangle$ be non-degenerate, and use intermediate normalization $\langle \Phi_i^{(0)} | \Psi_i \rangle = 1$. Then
\begin{align*}
E_i^{(1)} &= \langle \Phi_i^{(0)} | \hat{V} | \Phi_i^{(0)} \rangle, \\
|\Psi_i^{(1)}\rangle &= \sum_{n \neq i} \frac{|\Phi_n^{(0)}\rangle \langle \Phi_n^{(0)}|\hat{V}|\Phi_i^{(0)}\rangle}{E_i^{(0)} - E_n^{(0)}}
\end{align*}
The second-order energy is $E_i^{(2)} = \langle \Phi_i^{(0)}|\hat{V}|\Psi_i^{(1)}\rangle$, i.e.
\begin{align}\label{eq:e2-rspt}
E_i^{(2)} &= \sum_{n \neq i} \frac{|\langle \Phi_n^{(0)} | \hat{V} | \Phi_i^{(0)} \rangle|^2}{E_i^{(0)} - E_n^{(0)}}
\end{align}

\section{The Møller--Plesset Partitioning}

Take $\hat{H}_0$ as the sum of Fock operators (canonical HF orbitals):
\begin{align*}
\hat{H}_0 &= \sum_i \hat{f}(i), \quad \hat{V} = \hat{H} - \hat{H}_0 = \sum_{i<j} r_{ij}^{-1} - \sum_i v^{\text{HF}}(i)
\end{align*}
where $v^{\mathrm{HF}} = \sum_j (\hat{J}_j - \hat{K}_j)$. Every determinant built from HF orbitals is an eigenstate of $\hat{H}_0$ with eigenvalue equal to the sum of its occupied orbital energies. Hence
\begin{align*}
E_0^{(0)} &= \sum_i \epsilon_i, \\
E_0^{(0)} + E_0^{(1)} &= \langle \Phi_0 | \hat{H}_0 + \hat{V} | \Phi_0 \rangle = \langle \Phi_0 | \hat{H} | \Phi_0 \rangle = E_{\text{HF}}
\end{align*}
The second-order energy \eqref{eq:e2-rspt} needs $\langle \Phi_n|\hat{V}|\Phi_0\rangle$: since the matrix elements of $\hat{H}_0$ between different determinants vanish (it is diagonal in determinants), $\langle \Phi_n|\hat{V}|\Phi_0\rangle = \langle \Phi_n|\hat{H}|\Phi_0\rangle$, and only double excitations contribute:
\begin{align*}
\langle \Phi_{ij}^{ab} | \hat{H} | \Phi_0 \rangle &= \asym{ab}{ij}, \\
E_{ij}^{ab,(0)} - E_0^{(0)} &= \epsilon_a + \epsilon_b - \epsilon_i - \epsilon_j
\end{align*}
Summing over $i < j$, $a < b$ (each doubly excited determinant counted once):
\begin{align}\label{eq:mp2-so}
E^{(2)} &= \sum_{i<j} \sum_{a<b} \frac{|\asym{ij}{ab}|^2}{\epsilon_i + \epsilon_j - \epsilon_a - \epsilon_b} \\
&= \frac{1}{4} \sum_{ijab} \frac{|\asym{ij}{ab}|^2}{\epsilon_i + \epsilon_j - \epsilon_a - \epsilon_b} \notag
\end{align}

\section{Closed-Shell MP2: Spin Integration}

Replace the spin orbitals of \eqref{eq:mp2-so} by spatial orbitals $i, j, a, b$ plus spins. Let the spatial integrals be $\chem{ia}{jb}$ (chemists' notation). There are four cases according to the spins of $(I, J)$ ($I = i \sigma_i$ etc.). From Chapter 2, $\phys{IJ}{AB} = \chem{IA}{JB} = \chem{ia}{jb} \delta_{\sigma_I \sigma_A} \delta_{\sigma_J \sigma_B}$.

\textbf{$(\alpha, \alpha)$:} $A, B$ must also be $\alpha$:
\begin{align*}
\asym{i\alpha j\alpha}{a\alpha b\alpha} &= \chem{ia}{jb} - \chem{ib}{ja}
\end{align*}
Contribution $\frac{1}{4} \sum_{ijab} \frac{[\chem{ia}{jb} - \chem{ib}{ja}]^2}{D_{ij}^{ab}}$, with $D_{ij}^{ab} = \epsilon_i + \epsilon_j - \epsilon_a - \epsilon_b$.

\textbf{$(\beta, \beta)$:} the same as above.

\textbf{$(\alpha, \beta)$:} $(A, B)$ can be $(a \alpha, b \beta)$ or $(a \beta, b \alpha)$:
\begin{align*}
\asym{i\alpha j\beta}{a\alpha b\beta} &= \chem{ia}{jb} - 0, \quad \asym{i\alpha j\beta}{a\beta b\alpha} = 0 - \chem{ib}{ja}
\end{align*}
After squaring and summing over $a, b$, the second becomes identical to the first upon swapping the dummy indices $a \leftrightarrow b$ ($D$ is symmetric). Contribution $\frac{1}{4} \cdot 2 \sum \frac{\chem{ia}{jb}^2}{D}$.

\textbf{$(\beta, \alpha)$:} the same as $(\alpha, \beta)$.

In total:
\begin{align*}
E^{(2)} &= \frac{1}{2} \sum_{ijab} \frac{[\chem{ia}{jb} - \chem{ib}{ja}]^2}{D} + \sum_{ijab} \frac{\chem{ia}{jb}^2}{D}
\end{align*}
Expand the first term and use $\sum \frac{\chem{ib}{ja}^2}{D} = \sum \frac{\chem{ia}{jb}^2}{D}$ (swap the dummy indices $a \leftrightarrow b$):
\begin{align*}
\frac{1}{2} \sum \frac{[\chem{ia}{jb}^2 - 2 \chem{ia}{jb}\chem{ib}{ja} + \chem{ib}{ja}^2]}{D} &= \sum \frac{[\chem{ia}{jb}^2 - \chem{ia}{jb}\chem{ib}{ja}]}{D}
\end{align*}
Adding the opposite-spin part:
\begin{align}\label{eq:rmp2}
E_{\text{RMP2}} &= \sum_{ijab} \frac{\chem{ia}{jb} [2 \chem{ia}{jb} - \chem{ib}{ja}]}{\epsilon_i + \epsilon_j - \epsilon_a - \epsilon_b}
\end{align}
It is often split into a same-spin part $E_{SS}$ and an opposite-spin part $E_{OS}$; the derivation above shows clearly where each part comes from.

Implementation: \code{mp2} in \code{hf\_to\_cc/mp2.py}.

\begin{codemap}
\begin{lstlisting}[frame=none, numbers=none, xleftmargin=0pt, backgroundcolor={}]
ovov = mo[:nocc, nocc:, :nocc, nocc:]
D = eo[:, None, None, None] - ev[None, :, None, None] + eo[None, None, :, None] - ev[None, None, None, :]
E_mp2 = np.sum(ovov * (2 * ovov - ovov.transpose(0, 3, 2, 1)) / D)
\end{lstlisting}
\code{ovov[i, a, j, b]} is $\chem{ia}{jb}$, and \code{ovov.transpose(0, 3, 2, 1)[i, a, j, b]} is \code{ovov[i, b, j, a]} $= \chem{ib}{ja}$; \code{D[i, a, j, b]} is $\epsilon_i + \epsilon_j - \epsilon_a - \epsilon_b$. The last line is \eqref{eq:rmp2}.
\end{codemap}

\begin{summary}
\begin{itemize}
\item RSPT: $E^{(1)} = \langle 0|V|0\rangle$, $E^{(2)} = \sum \frac{|\langle n|V|0\rangle|^2}{E_0 - E_n}$; intermediate normalization.
\item MP partitioning: $E^{(0)} + E^{(1)} = E_{\text{HF}}$; $E^{(2)} = \frac{1}{4} \sum \frac{|\asym{ij}{ab}|^2}{D}$.
\item Spin integration: $E_{\text{RMP2}} = \sum \frac{(ia|jb)[2(ia|jb) - (ib|ja)]}{D}$.
\end{itemize}
\end{summary}

\chapter{Configuration Interaction: Structure and Size Consistency}
\label{ch:ci}

CI is conceptually the simplest correlation method: expand the wave function as a linear combination of determinants and apply linear variation (Chapter 1). This chapter discusses the structure of CI and its greatest weakness---truncated CI is not size consistent.

\textbf{Parallel reading.} “qc” pp.~71--77 (classification of CI, size consistency) and “Adv qc” pp.~33--35 (correlation energy and size consistency: the CID correlation energy of $N$ independent H$_2$ molecules is proportional to $\sqrt{N}$, the same as ``An Exactly Solvable Model'' in this chapter). The model comes from Szabo--Ostlund §4.6 (truncated CI and the size-consistency problem); Helgaker §11.3 also uses it to discuss the CID energy; Shavitt--Bartlett §1.5--1.8 discuss the same problem from the viewpoint of ``disconnected clusters''.

\section{The CI Expansion}

With the HF determinant as reference, classify by excitation level:
\begin{align*}
|\Psi\rangle &= c_0 |\Phi_0\rangle + \sum_{ia} c_i^a |\Phi_i^a\rangle \\
&\quad + \sum_{i<j, a<b} c_{ij}^{ab} |\Phi_{ij}^{ab}\rangle \\
&\quad + \sum_{i<j<k, a<b<c} c_{ijk}^{abc} |\Phi_{ijk}^{abc}\rangle + \dots
\end{align*}
Including all possible determinants gives FCI, the exact solution of the Schrödinger equation within the given basis. For the same orbitals the trial spaces are nested, so the variational principle (Chapter 1) gives $E_{\text{exact}} \le E_{\text{FCI}} \le E_{\text{CISD}} \le E_{\text{HF}}$ (with CISD built on this HF determinant).

\section{Size Consistency}

\subsection{Definition}

For two non-interacting systems A and B, the exact energy satisfies $E(A + B) = E(A) + E(B)$. A method that satisfies this when A and B are infinitely far apart is called size consistent. The premise is that the reference of A+B is the product of the references of A and B, and that the variational space or truncation of the method is compatible with A and B. Under this premise HF, FCI, CC and MBPT are size consistent; truncated CI is not. An example where the premise fails: when H$_2$ dissociates into two H atoms, the product of the two open-shell H atoms is not in the variational space of closed-shell RHF, so RHF is not size consistent.

\subsection{An Exactly Solvable Model}

Take $M$ H$_2$ molecules infinitely far apart, each in a minimal basis (one bonding orbital $\sigma_g$, one antibonding orbital $\sigma_u$). For a single H$_2$: reference $|0\rangle = |\sigma_g^2\rangle$ and a single double excitation $|D\rangle = |\sigma_u^2\rangle$ (single excitations do not couple to the ground state by symmetry). For $M$ molecules the CID space contains the reference and the $M$ doubly excited determinants $|D_A\rangle$ (only molecule A excited) but not the quadruples $|D_A D_B\rangle$, and its correlation energy grows only like $\sqrt{M}$: the correlation energy per molecule tends to zero as $M \rightarrow \infty$. This is what ``truncated CI is not size extensive'' means.

\textbf{What is missing?} The FCI wave function is the product of the wave functions of the molecules:
\begin{align*}
\Psi &= \prod_A (|0_A\rangle + c |D_A\rangle) = |0\rangle + c \sum_A |D_A\rangle + c^2 \sum_{A<B} |D_A D_B\rangle + \dots
\end{align*}
The coefficients of the quadruples $|D_A D_B\rangle$ are exactly products of double-excitation coefficients, $c^2$---such ``disconnected'' higher excitations are precisely what truncated CI throws away. The exponential ansatz $e^{T_2}$ of CC automatically contains $\frac{1}{2} T_2^2$, which is why CC is size extensive (Chapter 11).

\begin{summary}
\begin{itemize}
\item Truncated CI is not size consistent: the CID correlation energy of $M$ independent H$_2$ molecules grows only like $\sqrt{M}$, and the correlation energy per molecule tends to zero as $M$ grows.
\item What is missing are ``disconnected'' higher excitations (quadruples $\sim c_D^2$).
\end{itemize}
\end{summary}

\chapter{Coupled-Cluster Theory and CCD}
\label{ch:cc}

Chapter 10 showed that what truncated CI lacks are ``disconnected'' higher excitations---quadruples in which two double excitations happen at once, whose coefficients are approximately products of double-excitation coefficients. Coupled cluster (CC) includes these product terms automatically through the exponential ansatz and therefore stays size extensive after truncation. This chapter follows the line of MEST pp.~33--46 (“CC” pp.~1--17 follow the same route: formal theory, energy equation, amplitude equations, Feynman diagrams; “qc” pp.~82--84 and “Adv qc” p.~35 explain why the exponential ansatz automatically contains disconnected terms such as $T_2 T_2$; for textbooks see Shavitt--Bartlett §1.8 and Ch.~9--10 and the Crawford--Schaefer review pp.~45--111): the exponential ansatz, the similarity transformation, the natural truncation of the BCH expansion, the energy equation and the CCD amplitude equations.

\section{The Exponential Ansatz}

\begin{align*}
|\Psi_{\text{CC}}\rangle &= e^{T} |\Phi_0\rangle, \quad T = T_1 + T_2 + T_3 + \dots
\end{align*}
\begin{align*}
T_1 &= \sum_{ia} t_i^a \cre{a} \ann{i}, \quad T_2 = \frac{1}{4} \sum_{ijab} t_{ij}^{ab} \cre{a} \cre{b} \ann{j} \ann{i} \\
T_n &= \left(\frac{1}{n!}\right)^2 \sum t_{ij\ldots}^{ab\ldots} \cre{a} \cre{b} \ldots \ann{j} \ann{i}
\end{align*}
The amplitudes are antisymmetric in the occupied indices and in the virtual indices separately ($t_{ij}^{ab} = -t_{ji}^{ab} = -t_{ij}^{ba}$). All excitation operators are products of ``quasiparticle creation operators'' ($\cre{a}$ creates a particle, $a_i$ creates a hole), and they commute with each other:
\begin{align*}
[T_m, T_n] &= 0
\end{align*}
For two single-excitation operators (MEST p.~34):
\begin{align*}
[\cre{a} \ann{i}, \cre{b} \ann{j}] &= \cre{a} \ann{i} \cre{b} \ann{j} - \cre{b} \ann{j} \cre{a} \ann{i} \\
\cre{b} \ann{j} \cre{a} \ann{i} &= -\cre{b} \cre{a} \ann{j} \ann{i} = \cre{a} \cre{b} \ann{j} \ann{i} = -\cre{a} \cre{b} \ann{i} \ann{j} = \cre{a} \ann{i} \cre{b} \ann{j} \eqnote{$i, j \neq a, b$} \\
[\cre{a} \ann{i}, \cre{b} \ann{j}] &= 0
\end{align*}
Each step swaps two neighboring operators whose anticommutator is zero ($\{\ann{j}, \cre{a}\} = \delta_{ja} = 0$, $\{\ann{i}, \cre{b}\} = \delta_{ib} = 0$), so it only changes the sign.

\textbf{Expansion up to double excitations} (MEST p.~33):
\begin{align*}
e^{T} |\Phi_0\rangle &= \left(1 + T + \frac{1}{2!} T^2 + \frac{1}{3!} T^3 + \dots\right) |\Phi_0\rangle \\
&= \left(1 + T_1 + T_2 + \frac{1}{2} (T_1^2 + T_1 T_2 + T_2 T_1 + T_2^2) + \frac{1}{6} T_1^3 + \dots\right) |\Phi_0\rangle \\
&= \left(1 + T_1 + T_2 + \frac{1}{2} T_1^2 + T_1 T_2 + \frac{1}{2} T_2^2 + \frac{1}{6} T_1^3 + \dots\right) |\Phi_0\rangle
\end{align*}
The last step uses $T_2 T_1 = T_1 T_2$: the mixed term appears twice in $\frac{1}{2} T^2$, the square terms once.
Compared with the CI coefficients, $c_{ij}^{ab} = t_{ij}^{ab} + (t_i^a t_j^b - t_i^b t_j^a)$, and the coefficients of quadruples contain $\frac{1}{2} T_2^2$---exactly what truncated CI is missing.

\textbf{Size extensivity.} For two non-interacting systems A and B, $T = T_A + T_B$ (the excitation operators act on their own orbitals and commute with each other), and
\begin{align*}
(1 + T_A)(1 + T_B) &= 1 + T_A + T_B + T_A T_B \\
e^{T_A} e^{T_B} &= e^{T_A + T_B} \eqnote{$[T_A, T_B] = 0$} \\
e^{T_A + T_B} |\Phi_0^A \Phi_0^B\rangle &= e^{T_A} |\Phi_0^A\rangle e^{T_B} |\Phi_0^B\rangle
\end{align*}
A product of fragment wave functions contains the product $T_A T_B$, which a sum of excitation operators (CI) lacks; the exponential of the sum $T_A + T_B$ is the product of the fragment exponentials, and for operators this requires $[T_A, T_B] = 0$ (“qc” pp.~82--83, “Adv qc” p.~35).
The wave function is a product and the energies are additive. With truncation (e.g.\ keeping only $T_1 + T_2$) one must also show that $T_A + T_B$ really solves the truncated equations. Since $H = H_A + H_B$ and $T_A$ commutes with $H_B$, $\bar{H} = \bar{H}_A + \bar{H}_B$. Projected onto determinants that excite only A, this gives the residual of A itself; projected onto determinants in which both A and B are excited, it vanishes, because $\bar{H}_A$ can only excite A and $\bar{H}_B$ only B. So the fragment solutions together solve the whole system, with energy $E = E_A + E_B$. What is truncated is the excitation level of the connected cluster operator $T$; products in $e^{T}$ still produce higher, disconnected excitations---for example, the product of a double excitation on each fragment is a quadruple excitation.

\section{Similarity Transformation and the Projected Equations}

Substitute $e^{T}|\Phi_0\rangle$ into $H |\Psi\rangle = E |\Psi\rangle$ and multiply from the left by $e^{-T}$:
\begin{align*}
e^{-T} H e^T |\Phi_0\rangle = E |\Phi_0\rangle, \quad \bar{H} = e^{-T} H e^T
\end{align*}
Project onto the reference and the excited determinants:
\begin{align}\label{eq:cc-eqs}
E = \langle \Phi_0|\bar{H}|\Phi_0\rangle, \quad 0 = \langle \Phi_\mu|\bar{H}|\Phi_0\rangle \quad (\mu = \text{single, double, ... excitations})
\end{align}
The similarity-transformed equations contain only connected terms (next section).

\section{The BCH Expansion and Its Natural Truncation}

\subsection{The Expansion}

The Baker--Campbell--Hausdorff expansion:
\begin{align*}
e^{-T} H e^{T} &= H + [H, T] + \frac{1}{2!} [[H, T], T] + \frac{1}{3!} [[[H, T], T], T] \\
&\quad + \frac{1}{4!} [[[[H, T], T], T], T] + \dots
\end{align*}

\subsection{Commutators Keep Only Connected Terms}

Use the normal-ordered Hamiltonian $\HN  = H - E_{\text{ref}} = \FN  + \VN$ (Chapter 7). $T$ itself is normal ordered (it contains only quasiparticle creation operators). By Wick's theorem,
\begin{align*}
\HN T &= \no{\HN T} + (\text{terms of } \HN T \text{ with at least one contraction})
\end{align*}
The operators of $T$ cannot be contracted with each other or with the quasiparticle creation operators of $\HN$, and with $T$ on the left and $\HN$ on the right, the quasiparticle \emph{creation} operators of $T$ stand to the left of the quasiparticle \emph{annihilation} operators of $\HN$---their contractions vanish. So $T \HN  = \no{T \HN } = \no{\HN  T}$ (reordering inside a normal-ordered product only changes the sign, and $T$ is a product of an even number of operators, so the sign does not change). Hence
\begin{align*}
[\HN , T] &= \HN  T - T \HN  = (\HN  T)_c
\end{align*}
where $(\dots)_c$ means that there is at least one contraction line between $\HN$ and $T$ (connected). Likewise, in nested commutators every $T$ must be connected to $\HN$ (the argument of MEST pp.~34--36):
\begin{align*}
[[\HN, T], T] &= ((\HN T) T)_c \quad (\text{both } T \text{ connected to } \HN)
\end{align*}
($T$ and $T$ cannot contract with each other.)

\subsection{Two Worked Examples}

The cancellation of the uncontracted terms can be followed explicitly with the simplest cluster operator $T_1$ and the Fock part $\FN$ of $\HN$ (MEST pp.~34--36, “CC” p.~6).

\textbf{Example 1: $[\FN, T_1]$.}
\begin{align*}
[\FN, T_1] &= \sum_{pq} f_{pq} \sum_{ia} t_i^a [\no{\cre{p} \ann{q}}, \no{\cre{a} \ann{i}}]
\end{align*}
In $\no{\cre{p} \ann{q}} \no{\cre{a} \ann{i}}$ the possible contractions are $\ann{q}$--$\cre{a}$ (particle line) and $\cre{p}$--$\ann{i}$ (hole line); in $\no{\cre{a} \ann{i}} \no{\cre{p} \ann{q}}$ there are none, because $\cre{a}$ and $\ann{i}$ are quasiparticle creation operators standing on the left:
\begin{align*}
[\no{\cre{p} \ann{q}}, \no{\cre{a} \ann{i}}] &= \no{\cre{p} \ann{q}} \no{\cre{a} \ann{i}} - \no{\cre{a} \ann{i}} \no{\cre{p} \ann{q}} \\
&= \no{\cre{p} \ann{q} \cre{a} \ann{i}} + \delta_{qa} \no{\cre{p} \ann{i}} + \delta_{pi} \no{\ann{q} \cre{a}} + \delta_{pi} \delta_{qa} \\
&\quad - \no{\cre{a} \ann{i} \cre{p} \ann{q}} \\
&= \delta_{qa} \no{\cre{p} \ann{i}} + \delta_{pi} \no{\ann{q} \cre{a}} + \delta_{pi} \delta_{qa}
\end{align*}
The two uncontracted strings cancel: inside the braces every swap of two operators changes the sign (MEST p.~34), and moving $\cre{p} \ann{q}$ to the front of $\no{\cre{a} \ann{i} \cre{p} \ann{q}}$ takes four swaps. Hence
\begin{align}
[\FN, T_1] &= \sum_{pia} f_{pa} t_i^a \no{\cre{p} \ann{i}} + \sum_{qia} f_{iq} t_i^a \no{\ann{q} \cre{a}} + \sum_{ia} f_{ia} t_i^a \label{eq:cc-fnt1}
\end{align}
The same result follows from the anticommutation relations alone (“CC” p.~2):
\begin{align*}
[\cre{p} \ann{q}, \cre{a} \ann{i}] &= \cre{p} \ann{q} \cre{a} \ann{i} - \cre{a} \ann{i} \cre{p} \ann{q} \\
&= \cre{p} (\delta_{aq} - \cre{a} \ann{q}) \ann{i} - \cre{a} (\delta_{ip} - \cre{p} \ann{i}) \ann{q} \\
&= \delta_{aq} \cre{p} \ann{i} - \delta_{ip} \cre{a} \ann{q} \eqnote{the two four-operator terms cancel}
\end{align*}
and with $\cre{p} \ann{i} = \no{\cre{p} \ann{i}} + \delta_{pi}$ (hole line) and $\cre{a} \ann{q} = \no{\cre{a} \ann{q}} = -\no{\ann{q} \cre{a}}$ (no contraction: $a$ is a particle index) this is again $\delta_{qa} \no{\cre{p} \ann{i}} + \delta_{pi} \no{\ann{q} \cre{a}} + \delta_{pi} \delta_{qa}$.

\textbf{Example 2: $\frac{1}{2} [[\FN, T_1], T_1]$.}
\begin{align*}
\frac{1}{2} [[\FN, T_1], T_1] &= \frac{1}{2} [\FN T_1 - T_1 \FN, T_1] \\
&= \frac{1}{2} \FN T_1^2 - T_1 \FN T_1 + \frac{1}{2} T_1^2 \FN
\end{align*}
The three products are expanded in turn with the generalized Wick theorem.

\emph{First product.}
\begin{align*}
\frac{1}{2} \FN T_1^2 &= \frac{1}{2} \sum_{pq} f_{pq} \sum_{ia} t_i^a \sum_{jb} t_j^b \no{\cre{p} \ann{q}} \no{\cre{a} \ann{i}} \no{\cre{b} \ann{j}}
\end{align*}
$\cre{p}$ can contract with $\ann{i}$ or $\ann{j}$ (hole lines), $\ann{q}$ with $\cre{a}$ or $\cre{b}$ (particle lines):
\begin{align*}
\no{\cre{p} \ann{q}} \no{\cre{a} \ann{i}} \no{\cre{b} \ann{j}} &= \no{\cre{p} \ann{q} \cre{a} \ann{i} \cre{b} \ann{j}} + \delta_{pi} \no{\ann{q} \cre{a} \cre{b} \ann{j}} + \delta_{pj} \no{\ann{q} \cre{a} \ann{i} \cre{b}} \\
&\quad + \delta_{qa} \no{\cre{p} \ann{i} \cre{b} \ann{j}} + \delta_{qb} \no{\cre{p} \cre{a} \ann{i} \ann{j}} \\
&\quad + \delta_{pi} \delta_{qa} \no{\cre{b} \ann{j}} - \delta_{pi} \delta_{qb} \no{\cre{a} \ann{j}} + \delta_{pj} \delta_{qa} \no{\ann{i} \cre{b}} + \delta_{pj} \delta_{qb} \no{\cre{a} \ann{i}}
\end{align*}
(The minus sign: bringing $\ann{i}$ next to $\cre{p}$ takes two swaps, then bringing $\cre{b}$ next to $\ann{q}$ takes one.) Multiplying by $f_{pq}$ and summing over $p, q$:
\begin{align*}
\frac{1}{2} \FN T_1^2 &= \frac{1}{2} \sum_{iajb} t_i^a t_j^b \sum_{pq} f_{pq} \no{\cre{p} \ann{q} \cre{a} \ann{i} \cre{b} \ann{j}} \\
&\quad + \frac{1}{2} \sum_{iajb} t_i^a t_j^b \Big(\sum_q f_{iq} \no{\ann{q} \cre{a} \cre{b} \ann{j}} + \sum_q f_{jq} \no{\ann{q} \cre{a} \ann{i} \cre{b}}\Big) \\
&\quad + \frac{1}{2} \sum_{iajb} t_i^a t_j^b \Big(\sum_p f_{pa} \no{\cre{p} \ann{i} \cre{b} \ann{j}} + \sum_p f_{pb} \no{\cre{p} \cre{a} \ann{i} \ann{j}}\Big) \\
&\quad + \frac{1}{2} \sum_{iajb} t_i^a t_j^b \big(f_{ia} \no{\cre{b} \ann{j}} + f_{jb} \no{\cre{a} \ann{i}}\big) \\
&\quad + \frac{1}{2} \sum_{iajb} t_i^a t_j^b \big({-f_{ib}} \no{\cre{a} \ann{j}} + f_{ja} \no{\ann{i} \cre{b}}\big)
\end{align*}
The two terms in each pair of brackets are equal: relabeling the dummy indices $i \leftrightarrow j$, $a \leftrightarrow b$ turns one into the other. For the first pair:
\begin{align*}
\sum_{iajb} t_i^a t_j^b \sum_q f_{iq} \no{\ann{q} \cre{a} \cre{b} \ann{j}} &= \sum_{iajb} t_j^b t_i^a \sum_q f_{jq} \no{\ann{q} \cre{b} \cre{a} \ann{i}} \eqnote{$i \leftrightarrow j$, $a \leftrightarrow b$} \\
&= \sum_{iajb} t_i^a t_j^b \sum_q f_{jq} \no{\ann{q} \cre{a} \ann{i} \cre{b}} \eqnote{two swaps}
\end{align*}
and likewise for the others (in the last pair $-f_{ja} \no{\cre{b} \ann{i}} = f_{ja} \no{\ann{i} \cre{b}}$). So
\begin{align}
\frac{1}{2} \FN T_1^2 &= \frac{1}{2} \sum_{iajb} t_i^a t_j^b \sum_{pq} f_{pq} \no{\cre{p} \ann{q} \cre{a} \ann{i} \cre{b} \ann{j}} \notag \\
&\quad + \sum_{iajb} t_i^a t_j^b \Big(\sum_q f_{jq} \no{\ann{q} \cre{a} \ann{i} \cre{b}} + \sum_p f_{pb} \no{\cre{p} \cre{a} \ann{i} \ann{j}}\Big) \notag \\
&\quad + \sum_{iajb} t_i^a t_j^b \big(f_{jb} \no{\cre{a} \ann{i}} + f_{ja} \no{\ann{i} \cre{b}}\big) \label{eq:cc-ft1sq}
\end{align}

\emph{Second product.}
\begin{align*}
T_1 \FN T_1 &= \sum_{pq} f_{pq} \sum_{ia} t_i^a \sum_{jb} t_j^b \no{\cre{a} \ann{i}} \no{\cre{p} \ann{q}} \no{\cre{b} \ann{j}}
\end{align*}
The left $T_1$ cannot be contracted (its operators are quasiparticle creators standing on the left), so $\FN$ contracts only with the right $T_1$:
\begin{align*}
\no{\cre{a} \ann{i}} \no{\cre{p} \ann{q}} \no{\cre{b} \ann{j}} &= \no{\cre{a} \ann{i} \cre{p} \ann{q} \cre{b} \ann{j}} + \delta_{pj} \no{\cre{a} \ann{i} \ann{q} \cre{b}} \\
&\quad + \delta_{qb} \no{\cre{a} \ann{i} \cre{p} \ann{j}} + \delta_{pj} \delta_{qb} \no{\cre{a} \ann{i}}
\end{align*}
Inside the braces the pair $\cre{a} \ann{i}$ moves past any operator with two swaps, so in the operator order of \eqref{eq:cc-ft1sq}:
\begin{align}
T_1 \FN T_1 &= \sum_{iajb} t_i^a t_j^b \sum_{pq} f_{pq} \no{\cre{p} \ann{q} \cre{a} \ann{i} \cre{b} \ann{j}} \notag \\
&\quad + \sum_{iajb} t_i^a t_j^b \Big(\sum_q f_{jq} \no{\ann{q} \cre{a} \ann{i} \cre{b}} + \sum_p f_{pb} \no{\cre{p} \cre{a} \ann{i} \ann{j}}\Big) \notag \\
&\quad + \sum_{iajb} t_i^a t_j^b f_{jb} \no{\cre{a} \ann{i}} \label{eq:cc-t1ft1}
\end{align}

\emph{Third product.} Both $T_1$'s stand on the left, so nothing can be contracted:
\begin{align}
\frac{1}{2} T_1^2 \FN &= \frac{1}{2} \sum_{iajb} t_i^a t_j^b \sum_{pq} f_{pq} \no{\cre{a} \ann{i} \cre{b} \ann{j} \cre{p} \ann{q}} \notag \\
&= \frac{1}{2} \sum_{iajb} t_i^a t_j^b \sum_{pq} f_{pq} \no{\cre{p} \ann{q} \cre{a} \ann{i} \cre{b} \ann{j}} \label{eq:cc-t1sqf}
\end{align}

\emph{Sum.}
\begin{align}
\frac{1}{2} [[\FN, T_1], T_1] &= \eqref{eq:cc-ft1sq} - \eqref{eq:cc-t1ft1} + \eqref{eq:cc-t1sqf} \notag \\
&= \sum_{iajb} t_i^a t_j^b f_{ja} \no{\ann{i} \cre{b}} \label{eq:cc-fnt1t1}
\end{align}
The six-operator strings cancel ($\frac{1}{2} - 1 + \frac{1}{2} = 0$), and so do the $f_{jq}$, $f_{pb}$ and $f_{jb}$ terms ($1 - 1 = 0$): these are the terms in which at least one $T_1$ is not contracted with $\FN$. In the surviving term $\FN$ is contracted with both $T_1$'s ($\cre{p}$ with $\ann{j}$, $\ann{q}$ with $\cre{a}$).

In both examples the uncontracted strings are eliminated and only contracted terms survive: each commutator with $T$ keeps only the terms in which the new $T$ is contracted with what stands to its left. Generalizing (MEST p.~36): the only non-zero terms of the BCH expansion are those in which $\HN$ has at least one contraction with every $T$ on its right---the connected terms.

\subsection{Truncation at Fourth Order}

$\HN$ contains at most 4 operators ($\VN$), and each $T$ must take up at least one operator of $\HN$ for its contraction, so at most 4 $T$'s can be connected to $\HN$ at the same time. Five-fold and higher commutators vanish identically:
\begin{align*}
\bar{H} &= \HN  + (\HN  T)_c + \frac{1}{2} (\HN  T^2)_c + \frac{1}{6} (\HN  T^3)_c + \frac{1}{24} (\HN  T^4)_c + E_{\text{ref}} \\
&= (\HN  e^T)_c + E_{\text{ref}}
\end{align*}
This is why the CC equations can be written exactly with a finite number of terms.

\section{The Energy Equation}

\begin{align*}
E_{\text{CC}} - E_{\text{ref}} &= \langle \Phi_0|(\HN  e^T)_c|\Phi_0\rangle
\end{align*}
For $\langle\Phi_0|\dots|\Phi_0\rangle$ only fully contracted terms are non-zero. $\HN$ has at most 4 operators, so only $T_1$ (2 operators), $T_2$ (4) and $\frac{1}{2} T_1^2$ (4) can be fully contracted:
\begin{align*}
E_{\text{CC}} - E_{\text{ref}} &= \langle \Phi_0|(\FN  T_1)_c|\Phi_0\rangle \\
&\quad + \langle \Phi_0|(\VN  T_2)_c|\Phi_0\rangle + \frac{1}{2} \langle \Phi_0|(\VN  T_1^2)_c|\Phi_0\rangle
\end{align*}
($\FN  T_2$ and $\VN  T_1$ have mismatched numbers of operators and cannot be fully contracted. “CC” pp.~6--8 compute these three terms one by one with Wick contractions.)

\textbf{First term.} $\FN  T_1 = \sum_{pq} f_{pq} \sum_{ia} t_i^a \no{\cre{p} \ann{q}} \no{\cre{a} \ann{i}}$. Full contraction requires $\cre{p}$ to contract with $a_i$ (hole line: order $\cre{i} \ann{i}$, needs $p = i$) and $\ann{q}$ with $\cre{a}$ (particle line, $q = a$):
\begin{align*}
\langle \Phi_0|(\FN  T_1)_c|\Phi_0\rangle &= \sum_{ia} f_{ia} t_i^a
\end{align*}
This is the constant term of \eqref{eq:cc-fnt1}: $(\FN T_1)_c = [\FN, T_1]$, and its normal-ordered terms vanish in $\langle \Phi_0|\dots|\Phi_0\rangle$ (MEST p.~37). With HF orbitals $f_{ia} = 0$ (Brillouin's theorem) and the first term vanishes.

\textbf{Second term.} $\VN  T_2 = \frac{1}{4} \sum \asym{pq}{rs} \no{\cre{p} \cre{q} \ann{s} \ann{r}} \cdot \frac{1}{4} \sum t_{ij}^{ab} \no{\cre{a} \cre{b} \ann{j} \ann{i}}$, so
\begin{align*}
\langle \Phi_0|(\VN  T_2)_c|\Phi_0\rangle &= \frac{1}{16} \sum_{pqrs} \asym{pq}{rs} \sum_{ijab} t_{ij}^{ab} \langle \Phi_0|\no{\cre{p} \cre{q} \ann{s} \ann{r}} \no{\cre{a} \cre{b} \ann{j} \ann{i}}|\Phi_0\rangle
\end{align*}
$\cre{p}, \cre{q}$ must contract with $\ann{j}, \ann{i}$ (hole lines) and $\ann{s}, \ann{r}$ with $\cre{a}, \cre{b}$ (particle lines), in four ways. With the operators numbered $\cre{p} (1)\, \cre{q} (2)\, \ann{s} (3)\, \ann{r} (4)\allowbreak\, |\, \cre{a} (5)\, \cre{b} (6)\, \ann{j} (7)\, \ann{i} (8)$, the four full contractions $(1,8)(2,7)(3,6)(4,5)$, $(1,7)(2,8)(3,6)(4,5)$, $(1,8)(2,7)(3,5)(4,6)$ and $(1,7)(2,8)(3,5)(4,6)$ have 0, 1, 1 and 2 crossings (MEST p.~37, “CC” pp.~7--8):
\begin{align*}
\langle \Phi_0|\no{\cre{p} \cre{q} \ann{s} \ann{r}} \no{\cre{a} \cre{b} \ann{j} \ann{i}}|\Phi_0\rangle &= \delta_{pi} \delta_{qj} \delta_{sb} \delta_{ra} - \delta_{pj} \delta_{qi} \delta_{sb} \delta_{ra} \\
&\quad - \delta_{pi} \delta_{qj} \delta_{sa} \delta_{rb} + \delta_{pj} \delta_{qi} \delta_{sa} \delta_{rb} \\
\langle \Phi_0|(\VN  T_2)_c|\Phi_0\rangle &= \frac{1}{16} \sum_{ijab} t_{ij}^{ab} (\asym{ij}{ab} - \asym{ji}{ab} - \asym{ij}{ba} + \asym{ji}{ba}) \\
&= \frac{1}{4} \sum_{ijab} \asym{ij}{ab} t_{ij}^{ab} \eqnote{$\asym{ji}{ab} = \asym{ij}{ba} = -\asym{ij}{ab}$}
\end{align*}
The four contractions are equal, and $4 \cdot \frac{1}{16} = \frac{1}{4}$.

\textbf{Third term.} $T_1^2 = \sum_{ia} \sum_{jb} t_i^a t_j^b \{\cre{a} \ann{i}\}\{\cre{b} \ann{j}\}$, so
\begin{align*}
\frac{1}{2} \langle \Phi_0|(\VN  T_1^2)_c|\Phi_0\rangle &= \frac{1}{8} \sum_{pqrs} \asym{pq}{rs} \sum_{iajb} t_i^a t_j^b \langle \Phi_0|\no{\cre{p} \cre{q} \ann{s} \ann{r}} \no{\cre{a} \ann{i}} \no{\cre{b} \ann{j}}|\Phi_0\rangle
\end{align*}
The $\cre{p}, \cre{q}$ of $\VN$ must contract with the two hole operators $\ann{i}, \ann{j}$, and $\ann{s}, \ann{r}$ with the two particle operators $\cre{a}, \cre{b}$, again in four ways. With the operators numbered $\cre{p} (1)\, \cre{q} (2)\, \ann{s} (3)\, \ann{r} (4)\allowbreak\, |\, \cre{a} (5)\, \ann{i} (6)\allowbreak\, |\, \cre{b} (7)\, \ann{j} (8)$, the four full contractions $(1,8)(2,6)(3,7)(4,5)$, $(1,8)(2,6)(3,5)(4,7)$, $(1,6)(2,8)(3,7)(4,5)$ and $(1,6)(2,8)(3,5)(4,7)$ have 1, 2, 2 and 3 crossings (MEST p.~37):
\begin{align*}
\langle \Phi_0|\no{\cre{p} \cre{q} \ann{s} \ann{r}} \no{\cre{a} \ann{i}} \no{\cre{b} \ann{j}}|\Phi_0\rangle &= -\delta_{pj} \delta_{qi} \delta_{sb} \delta_{ra} + \delta_{pj} \delta_{qi} \delta_{sa} \delta_{rb} \\
&\quad + \delta_{pi} \delta_{qj} \delta_{sb} \delta_{ra} - \delta_{pi} \delta_{qj} \delta_{sa} \delta_{rb} \\
\frac{1}{2} \langle \Phi_0|(\VN  T_1^2)_c|\Phi_0\rangle &= \frac{1}{8} \sum_{iajb} t_i^a t_j^b (-\asym{ji}{ab} + \asym{ji}{ba} + \asym{ij}{ab} - \asym{ij}{ba}) \\
&= \frac{1}{2} \sum_{ijab} \asym{ij}{ab} t_i^a t_j^b \eqnote{$4 \cdot \frac{1}{8} = \frac{1}{2}$}
\end{align*}
Altogether:
\begin{align}\label{eq:ecc}
E_{\text{CC}} &= E_{\text{ref}} + \sum_{ia} f_{ia} t_i^a + \frac{1}{4} \sum_{ijab} \asym{ij}{ab} t_{ij}^{ab} \\
&\quad + \frac{1}{2} \sum_{ijab} \asym{ij}{ab} t_i^a t_j^b \notag
\end{align}
This formula holds for CCSD, CCSDT, \ldots\ (higher $T_n$ cannot be fully contracted with $\HN$).

\section{The CCD Amplitude Equations}

Take $T = T_2$:
\begin{align*}
0 &= \langle \Phi_{ij}^{ab}|(\HN  + \HN  T_2 + \frac{1}{2} \HN  T_2^2)_c|\Phi_0\rangle
\end{align*}
($T_2^3$ has 12 operators; after connecting to $\HN$ at least 8 remain, which cannot be fully contracted with the 4 of $\langle \Phi_{ij}^{ab}|$.)

\subsection{Zeroth-Order and Fock Terms}

$\langle \Phi_{ij}^{ab}|\VN |\Phi_0\rangle = \asym{ab}{ij}$ (Chapter 7). The Fock term $\langle \Phi_{ij}^{ab}|(\FN  T_2)_c|\Phi_0\rangle$: one operator of $\FN$ contracts with $T_2$ and the other with $\langle \Phi_{ij}^{ab}| = \langle \Phi_0|\{\cre{i} \cre{j} \ann{b} \ann{a}\}$ on the left. Two classes: particle lines ($f_{bc}$ replaces the $c$ of $T_2$ by $b$) and hole lines ($f_{kj}$). Result:
\begin{align*}
\langle \Phi_{ij}^{ab}|(\FN  T_2)_c|\Phi_0\rangle &= P(ab) \sum_c f_{bc} t_{ij}^{ac} - P(ij) \sum_k f_{kj} t_{ik}^{ab}
\end{align*}
Write out the sign for one particle-line term. Number the operators 1--10 (the four operators of $T_2$ written last):
\begin{align*}
\cre{i} (1)  \cre{j} (2)  \ann{b} (3)  \ann{a} (4)  |  \cre{p} (5)  \ann{q} (6)  |  \cre{c} (7)  \cre{d} (8)  \ann{l} (9)  \ann{k} (10)
\end{align*}
Take $\ann{b}$--$\cre{p}$ ($p = b$), $\ann{q}$--$\cre{c}$ ($q = c$), $\ann{a}$--$\cre{d}$ ($d = a$), $\cre{j}$--$a_l$ ($l = j$), $\cre{i}$--$a_k$ ($k = i$). Of the contraction pairs $(1,10), (2,9), (3,5), (4,8), (6,7)$ only $(3,5)$ and $(4,8)$ cross, sign $-$, value $-f_{bc} t_{ij}^{ca} = +f_{bc} t_{ij}^{ac}$. If instead $\ann{a}$--$\cre{p}$ and $\ann{b}$--$\cre{d}$, the contraction lines $(4,5), (3,8), (6,7), (2,9), (1,10)$ no longer cross, and the value is $f_{ac} t_{ij}^{cb} = -f_{ac} t_{ij}^{bc}$. Together they make $P(ab) \sum_c f_{bc} t_{ij}^{ac}$: the external indices $a, b$ can come from $\FN$ or from $T_2$, and the two possibilities differ by a sign.

Here $P(ab) = 1 - P_{ab}$, where $P_{ab}$ interchanges the labels $a$ and $b$ in the expression it acts on; likewise $P(ij) = 1 - P_{ij}$, and $P(ij) P(ab)$ applies both (“CC” p.~17):
\begin{align*}
P(ab) \sum_c f_{bc} t_{ij}^{ac} &= \sum_c f_{bc} t_{ij}^{ac} - \sum_c f_{ac} t_{ij}^{bc} \\
-P(ij) \sum_k f_{kj} t_{ik}^{ab} &= -\sum_k f_{kj} t_{ik}^{ab} + \sum_k f_{ki} t_{jk}^{ab}
\end{align*}
An antisymmetrizer is needed only for a pair of external indices that sit on different factors ($a$ on $t$ and $b$ on $f$ in $f_{bc} t_{ij}^{ac}$); a pair that sits on one antisymmetric factor ($i, j$ in $t_{ij}^{ac}$) is antisymmetric already (“CC” pp.~16--17, MEST p.~43).

\subsection{The Complete CCD Equations}

Adding the three linear terms of $(\VN T_2)_c$ (the particle--particle and hole--hole ladder terms and the ring term) and the four quadratic terms of $\frac{1}{2} (\VN T_2^2)_c$ gives the complete CCD equations:
\begin{align}
0 &= \asym{ab}{ij} + P(ab) \sum_c f_{bc} t_{ij}^{ac} - P(ij) \sum_k f_{kj} t_{ik}^{ab} \notag \\
&\quad + \frac{1}{2} \sum_{cd} \asym{ab}{cd} t_{ij}^{cd} + \frac{1}{2} \sum_{kl} \asym{kl}{ij} t_{kl}^{ab} + P(ij) P(ab) \sum_{kc} \asym{kb}{cj} t_{ik}^{ac} \notag \\
&\quad + \frac{1}{4} \sum_{klcd} \asym{kl}{cd} t_{ij}^{cd} t_{kl}^{ab} + P(ij) \sum_{klcd} \asym{kl}{cd} t_{ik}^{ac} t_{jl}^{bd} \notag \\
&\quad - \frac{1}{2} P(ab) \sum_{klcd} \asym{kl}{cd} t_{ij}^{ac} t_{kl}^{bd} - \frac{1}{2} P(ij) \sum_{klcd} \asym{kl}{cd} t_{ik}^{ab} t_{jl}^{cd} \label{eq:ccd}
\end{align}
This agrees with Shavitt--Bartlett Eq.~(9.126) (p.~288 of the book, derived with diagrams, §9.5) and with Eq.~(15.10) of Python for QC (which factorizes the quadratic terms into intermediates).

The residual $R_{ij}^{ab} = \langle \Phi_{ij}^{ab}|\bar{H}|\Phi_0\rangle$ vanishes at convergence. The iteration starts from $t_{ij}^{ab} = \asym{ij}{ab} / D_{ij}^{ab}$, exactly the first-order amplitude (inserting it into the energy formula gives the MP2 energy), and updates $t \leftarrow t + R / D$ with $D_{ij}^{ab} = f_{ii} + f_{jj} - f_{aa} - f_{bb}$.

Implementation: \code{ccd} in \code{hf\_to\_cc/ccd.py}.

\begin{codemap}
\begin{lstlisting}[frame=none, numbers=none, xleftmargin=0pt, backgroundcolor={}]
R = oovv + quad(t_o)
R += 0.5 * np.einsum("abcd, ijcd->ijab", vvvv, t_o)
R += 0.5 * np.einsum("klij, klab->ijab", oooo, t_o)
X = np.einsum("kbcj, ikac->ijab", ovvo, t_o)
R += X - X.transpose(1, 0, 2, 3) - X.transpose(0, 1, 3, 2) + X.transpose(1, 0, 3, 2)
t_n = R / D
\end{lstlisting}
Line by line, \eqref{eq:ccd}: $\asym{ab}{ij}$ (\code{oovv}, equal to $\asym{ij}{ab}$ for real orbitals) and the four quadratic terms (\code{quad}, called from \code{provided.py}); the two ladder terms; the ring term $X_{ij}^{ab} = \sum_{kc} \asym{kb}{cj} t_{ik}^{ac}$, whose three transposes are $P(ij)$, $P(ab)$ and $P(ij) P(ab)$ with their signs. The two Fock terms are absent: for canonical orbitals they equal $-D_{ij}^{ab} t_{ij}^{ab}$, so moving them to the other side gives $t = R / D$ with $R$ the remaining terms.
\end{codemap}

\begin{summary}
\begin{itemize}
\item $\bar{H} = e^{-T} H e^T = (\HN  e^T)_c + E_{\text{ref}}$; the BCH expansion terminates naturally at the fourfold commutator.
\item The energy depends only on $T_1$ and $T_2$: $E = \sum f_{ia} t_{ia} + \frac{1}{4} \sum \asym{ij}{ab} t_{ij}^{ab} + \frac{1}{2} \sum \asym{ij}{ab} t_i^a t_j^b$.
\item CCD equations = zeroth order + Fock + three linear terms (two ladders, one ring) + four quadratic terms.
\end{itemize}
\end{summary}


\end{document}
